% Options for packages loaded elsewhere
\PassOptionsToPackage{unicode}{hyperref}
\PassOptionsToPackage{hyphens}{url}
\PassOptionsToPackage{dvipsnames,svgnames,x11names}{xcolor}
%
\documentclass[
  10pt,
]{article}
\usepackage{amsmath,amssymb}
\usepackage{iftex}
\ifPDFTeX
  \usepackage[T1]{fontenc}
  \usepackage[utf8]{inputenc}
  \usepackage{textcomp} % provide euro and other symbols
\else % if luatex or xetex
  \usepackage{unicode-math} % this also loads fontspec
  \defaultfontfeatures{Scale=MatchLowercase}
  \defaultfontfeatures[\rmfamily]{Ligatures=TeX,Scale=1}
\fi
\usepackage{lmodern}
\ifPDFTeX\else
  % xetex/luatex font selection
  \setmainfont[]{DejaVu Serif}
  \setsansfont[]{DejaVu Sans}
  \setmonofont[]{DejaVu Sans Mono}
  \setmathfont[]{Latin Modern Math}
\fi
% Use upquote if available, for straight quotes in verbatim environments
\IfFileExists{upquote.sty}{\usepackage{upquote}}{}
\IfFileExists{microtype.sty}{% use microtype if available
  \usepackage[]{microtype}
  \UseMicrotypeSet[protrusion]{basicmath} % disable protrusion for tt fonts
}{}
\makeatletter
\@ifundefined{KOMAClassName}{% if non-KOMA class
  \IfFileExists{parskip.sty}{%
    \usepackage{parskip}
  }{% else
    \setlength{\parindent}{0pt}
    \setlength{\parskip}{6pt plus 2pt minus 1pt}}
}{% if KOMA class
  \KOMAoptions{parskip=half}}
\makeatother
\usepackage{xcolor}
\usepackage[margin=2.0cm]{geometry}
\usepackage{color}
\usepackage{fancyvrb}
\newcommand{\VerbBar}{|}
\newcommand{\VERB}{\Verb[commandchars=\\\{\}]}
\DefineVerbatimEnvironment{Highlighting}{Verbatim}{commandchars=\\\{\}}
% Add ',fontsize=\small' for more characters per line
\newenvironment{Shaded}{}{}
\newcommand{\AlertTok}[1]{\textcolor[rgb]{1.00,0.00,0.00}{\textbf{#1}}}
\newcommand{\AnnotationTok}[1]{\textcolor[rgb]{0.38,0.63,0.69}{\textbf{\textit{#1}}}}
\newcommand{\AttributeTok}[1]{\textcolor[rgb]{0.49,0.56,0.16}{#1}}
\newcommand{\BaseNTok}[1]{\textcolor[rgb]{0.25,0.63,0.44}{#1}}
\newcommand{\BuiltInTok}[1]{\textcolor[rgb]{0.00,0.50,0.00}{#1}}
\newcommand{\CharTok}[1]{\textcolor[rgb]{0.25,0.44,0.63}{#1}}
\newcommand{\CommentTok}[1]{\textcolor[rgb]{0.38,0.63,0.69}{\textit{#1}}}
\newcommand{\CommentVarTok}[1]{\textcolor[rgb]{0.38,0.63,0.69}{\textbf{\textit{#1}}}}
\newcommand{\ConstantTok}[1]{\textcolor[rgb]{0.53,0.00,0.00}{#1}}
\newcommand{\ControlFlowTok}[1]{\textcolor[rgb]{0.00,0.44,0.13}{\textbf{#1}}}
\newcommand{\DataTypeTok}[1]{\textcolor[rgb]{0.56,0.13,0.00}{#1}}
\newcommand{\DecValTok}[1]{\textcolor[rgb]{0.25,0.63,0.44}{#1}}
\newcommand{\DocumentationTok}[1]{\textcolor[rgb]{0.73,0.13,0.13}{\textit{#1}}}
\newcommand{\ErrorTok}[1]{\textcolor[rgb]{1.00,0.00,0.00}{\textbf{#1}}}
\newcommand{\ExtensionTok}[1]{#1}
\newcommand{\FloatTok}[1]{\textcolor[rgb]{0.25,0.63,0.44}{#1}}
\newcommand{\FunctionTok}[1]{\textcolor[rgb]{0.02,0.16,0.49}{#1}}
\newcommand{\ImportTok}[1]{\textcolor[rgb]{0.00,0.50,0.00}{\textbf{#1}}}
\newcommand{\InformationTok}[1]{\textcolor[rgb]{0.38,0.63,0.69}{\textbf{\textit{#1}}}}
\newcommand{\KeywordTok}[1]{\textcolor[rgb]{0.00,0.44,0.13}{\textbf{#1}}}
\newcommand{\NormalTok}[1]{#1}
\newcommand{\OperatorTok}[1]{\textcolor[rgb]{0.40,0.40,0.40}{#1}}
\newcommand{\OtherTok}[1]{\textcolor[rgb]{0.00,0.44,0.13}{#1}}
\newcommand{\PreprocessorTok}[1]{\textcolor[rgb]{0.74,0.48,0.00}{#1}}
\newcommand{\RegionMarkerTok}[1]{#1}
\newcommand{\SpecialCharTok}[1]{\textcolor[rgb]{0.25,0.44,0.63}{#1}}
\newcommand{\SpecialStringTok}[1]{\textcolor[rgb]{0.73,0.40,0.53}{#1}}
\newcommand{\StringTok}[1]{\textcolor[rgb]{0.25,0.44,0.63}{#1}}
\newcommand{\VariableTok}[1]{\textcolor[rgb]{0.10,0.09,0.49}{#1}}
\newcommand{\VerbatimStringTok}[1]{\textcolor[rgb]{0.25,0.44,0.63}{#1}}
\newcommand{\WarningTok}[1]{\textcolor[rgb]{0.38,0.63,0.69}{\textbf{\textit{#1}}}}
\usepackage{longtable,booktabs,array}
\usepackage{calc} % for calculating minipage widths
% Correct order of tables after \paragraph or \subparagraph
\usepackage{etoolbox}
\makeatletter
\patchcmd\longtable{\par}{\if@noskipsec\mbox{}\fi\par}{}{}
\makeatother
% Allow footnotes in longtable head/foot
\IfFileExists{footnotehyper.sty}{\usepackage{footnotehyper}}{\usepackage{footnote}}
\makesavenoteenv{longtable}
\setlength{\emergencystretch}{3em} % prevent overfull lines
\providecommand{\tightlist}{%
  \setlength{\itemsep}{0pt}\setlength{\parskip}{0pt}}
\setcounter{secnumdepth}{-\maxdimen} % remove section numbering
% doc_header.tex -- LaTeX preamble additions used by build_docs.sh for every
% Markdown -> LaTeX -> PDF conversion (the documents of the folder Revision_SM_Higgs).
\usepackage{etoolbox}
% longtable may not be loaded by pandoc when a document has no table
\usepackage{longtable}
% bold vectors: \bm -> \symbf under unicode-math (XeLaTeX), bm package otherwise
\makeatletter
\@ifpackageloaded{unicode-math}{\providecommand{\bm}[1]{\symbf{#1}}}{\usepackage{bm}}
\makeatother
% smaller type in long tables so that wide audit tables fit the text width;
% inside a table a line break is allowed after an underscore of plain text
% (pandoc writes it as \_), so that long identifiers wrap inside their column
% rows are set a little apart, so that symbols with indices do not touch
\AtBeginEnvironment{longtable}{\footnotesize\renewcommand{\arraystretch}{1.45}%
  \let\sptunderscore\_\def\_{\sptunderscore\allowbreak}}
\setlength{\LTcapwidth}{\textwidth}
% A short table is kept on one page.  The build (code/scripts/keep_tables.py) ends the rows of such a table with \\*,
% which forbids a page break after a row, and sets \sptkeptpenalty to 10000 in front of the table.  This penalty is
% put here in the place of the breakpoint (penalty 0) that longtable leaves after every table, so that the paragraph
% that follows a kept table, its caption, begins on the same page; it is set back to 0 after the table.
\newcount\sptkeptpenalty
\sptkeptpenalty=0
\makeatletter
\patchcmd{\endlongtable}{\endgraf\penalty\z@\addvspace\LTpost}%
  {\endgraf\penalty\sptkeptpenalty\addvspace\LTpost\global\sptkeptpenalty=\z@}%
  {}{\PackageError{doc_header}{the end of longtable could not be patched}{The tables can still be typeset, but a short table may be separated from its caption.}}
\makeatother
% allow line breaks after slashes and hyphens inside table cells
\usepackage{xurl}
% code blocks: break over-long lines instead of running into the margin.
% Plain ``` blocks are typeset by pandoc as {verbatim}, highlighted ones
% (```bash) as {Highlighting}; both are re-declared as breakable fancyvrb
% environments.  The continuation line is marked with a small hooked arrow.
\usepackage{fvextra}
\RecustomVerbatimEnvironment{verbatim}{Verbatim}{breaklines,breakanywhere,fontsize=\small}
\makeatletter
\@ifundefined{Highlighting}{}{%
  \RecustomVerbatimEnvironment{Highlighting}{Verbatim}{breaklines,breakanywhere,fontsize=\small,commandchars=\\\{\}}}
\makeatother
% inline code (pandoc writes it as \texttt{...}): never hyphenate -- an inserted
% hyphen could be read as part of a file name -- but allow a line break after
% an underscore, slash, comma, colon or equals sign, so that long paths and
% column lists wrap inside the text width and inside table columns.
% Inside a formula (e.g. the name of a notebook symbol in an equation) the same
% command must first leave math mode, which \text of amsmath does.
\ExplSyntaxOn
\tl_new:N \l__spt_code_tl
\cs_new_protected:Npn \__spt_texttt:n #1
  {
    \tl_set:Nn \l__spt_code_tl {#1}
    \regex_replace_all:nnN { ([\/\,\:\=]) } { \1 \c{allowbreak} } \l__spt_code_tl
    \regex_replace_all:nnN { \c{_} } { \c{_} \c{allowbreak} } \l__spt_code_tl
    \group_begin: \ttfamily \hyphenchar \font = -1 \scan_stop: \l__spt_code_tl \group_end:
  }
\RenewDocumentCommand \texttt { m }
  {
    \mode_if_math:TF { \text { \__spt_texttt:n {#1} } } { \__spt_texttt:n {#1} }
  }
\ExplSyntaxOff
% A figure is written in the Markdown source as an image paragraph followed by
% a caption paragraph ("**Figure S1.** ..."), so that the caption is also
% visible on GitHub.  An image that starts a paragraph is centered.
\usepackage{graphicx}% pandoc loads it only for documents that contain an image
\NewCommandCopy\sptincludegraphics\includegraphics
\RenewDocumentCommand\includegraphics{O{} m}{%
  \ifvmode\begin{center}\sptincludegraphics[#1]{#2}\end{center}%
  \else\sptincludegraphics[#1]{#2}\fi}
% keep figures where they are referenced
\usepackage{float}
\floatplacement{figure}{H}
% running header with the document title
\usepackage{fancyhdr}
\pagestyle{fancy}
\fancyhf{}
\fancyhead[L]{\small PROJECT spacetimes, Revision\_SM\_Higgs}
\fancyhead[R]{\small\thepage}
\renewcommand{\headrulewidth}{0.4pt}
% 16x16 matrices: tighter columns so that a full matrix fits the text width
\setlength{\arraycolsep}{3.2pt}
% The documents of this project contain long tagged symbols, file paths and
% identifiers that cannot be hyphenated.  A paragraph is therefore allowed to
% have looser lines before TeX gives up and lets a line stick out of the margin.
\AtBeginDocument{\tolerance=2000 \emergencystretch=5em \hbadness=2000}
\ifLuaTeX
  \usepackage{selnolig}  % disable illegal ligatures
\fi
\IfFileExists{bookmark.sty}{\usepackage{bookmark}}{\usepackage{hyperref}}
\IfFileExists{xurl.sty}{\usepackage{xurl}}{} % add URL line breaks if available
\urlstyle{same}
\hypersetup{
  pdftitle={The scale, the masses of the gauge bosons and the amplitude mode of the Higgs superfluid},
  pdfauthor={Prepared for Patrick Lee Nash, Ph.D.},
  colorlinks=true,
  linkcolor={blue},
  filecolor={Maroon},
  citecolor={blue},
  urlcolor={blue},
  pdfcreator={LaTeX via pandoc}}

\title{The scale, the masses of the gauge bosons and the amplitude mode
of the Higgs superfluid}
\usepackage{etoolbox}
\makeatletter
\providecommand{\subtitle}[1]{% add subtitle to \maketitle
  \apptocmd{\@title}{\par {\large #1 \par}}{}{}
}
\makeatother
\subtitle{Stage 4 of the revision of the hypotheses, second version,
folder Revision\_SM\_Higgs of PROJECT spacetimes}
\author{Prepared for Patrick Lee Nash, Ph.D.}
\date{November 21, 2026}

\begin{document}
\maketitle

\hypertarget{what-this-document-is-and-what-it-finds}{%
\section{0. What this document is and what it
finds}\label{what-this-document-is-and-what-it-finds}}

The author asked that the hypotheses of PROJECT spacetimes be modified
so that the superfluid can be the Higgs field of the Standard Model,
with the quantum numbers of the Higgs doublet. In the second version of
the revised hypotheses, stated in Section 1, the superfluid is a
condensate of pairs of a particle of the isodoublet \(\Psi_{\rm LL}\)
with a particle of one of the two singlets \(\chi_{\rm SH}\) and
\(\chi'_{\rm SH}\). Stage 3 computed its ground state in three
mean-field approximations, the methods of Hubbard and Stratonovich (HS,
with three parameter sets), of Kohn and Sham (KS) and of Cornwall,
Jackiw and Tomboulis (CJT). In every method the ground state is the
custodial Higgs superfluid: both pair fields condense with equal
lengths, and the state keeps the electric charge and the custodial
symmetry. This document is stage 4: it computes what the Standard Model
measures of its Higgs field. Those quantities are the scale of the
breaking, the masses of the weak bosons, the parameter \(\rho\), and the
mass of the amplitude mode, the particle that would be the Higgs boson.
The first version of the hypotheses (one singlet) is computed alongside
as the comparison. All quantities are first computed in units of the
momentum cutoff \(\Lambda_{\rm LL}=1\) and then converted to GeV by the
value of the scale, \(v_{\rm SH}=246.22\) GeV.

The findings:

\begin{enumerate}
\def\labelenumi{\arabic{enumi}.}
\tightlist
\item
  \textbf{The superfluid is a minimum of its effective potential.} Along
  the direction of the vacuum, the energy of the lowest state with a
  given pair amplitude is a Mexican hat. It is zero in the normal
  vacuum, it falls to the ground state of stage 3 and rises beyond it,
  in every method and both versions (Section 2).
\item
  \textbf{The photon is massless, and in the second version \(\rho=1\)
  exactly.} The susceptibilities of the superfluid to the chemical
  potentials of the generators are computed (Section 3). The electric
  charge has zero susceptibility in every method, so the mass matrix of
  the electroweak gauge bosons has one zero eigenvalue with the
  eigenvector of the photon. The second version has the custodial
  structure: in every method \(\rho_{\rm SH}\) is \(1\). The first
  version has \(\rho_{\rm SH}\) from \(1.202\) to \(1.247\), which the
  measured \(\rho\) excludes.
\item
  \textbf{The number of broken generators is four in the second
  version.} The susceptibility matrix of the eight generators has rank
  four in every method. The three generators of \(SU(2)_L\times U(1)_Y\)
  that the photon does not keep give the longitudinal parts of
  \(W^\pm_\mu\) and \(Z_\mu\); the fourth, the particle number, gives a
  massless Goldstone boson. The linear response at zero frequency finds
  the same four modes (Section 5).
\item
  \textbf{The scale.} \(v_{\rm SH}=2(\chi^{11}_{\rm SH})^{1/2}\) is
  \(0.102\) to \(0.1116\) in units of the cutoff in the second version.
  With \(v_{\rm SH}=246.22\) GeV the cutoff \(\Lambda_{\rm LL}\) is
  \(2207\) to \(2413\) GeV, and the lightest fermionic excitations,
  which carry the fractional electric charges \(\tfrac34\) and
  \(-\tfrac14\), have masses of \(848\) to \(1.09\times10^{3}\) GeV
  (Section 4).
\item
  \textbf{The amplitude mode lies exactly at the pair self-energy, and
  it is far too heavy.} In the pairing-only methods HS (0, 1) and HS (0,
  0) the linear response has a closed form, that of the theory of
  Bardeen, Cooper and Schrieffer for a superconductor. It puts the
  amplitude mode exactly at \(\omega^{\rm H}_{\rm SH}=d^0_{\rm SH}\),
  below the threshold of two quasi-particles, for every coupling and
  mass (Section 5.3). With the scale of the second version this is
  \(M_{\rm H}/v_{\rm SH}=7.967\), that is \(1962\) GeV. The effective
  Lagrangian gives \(9.22\) to \(9.57\) in all five methods. The
  Standard Model has \(m_{\rm H}/v=0.5082\): the superfluid's amplitude
  mode is \(16\) to \(19\) times too heavy (Section 5).
\item
  \textbf{No coupling repairs it.} In HS (0, 0), where the ground state
  has a closed form, \(M_{\rm H}/v_{\rm SH}\) lies between \(4.892\) and
  \(34.43\) for \(1.02\le G_{\rm SH}/G^{\rm crit}_{\rm SH}\le3\) and the
  masses 0.1 and 0.2 of the fields (Section 6). The mean-field
  superfluid of these hypotheses therefore does not have the mass of the
  Higgs boson.
\item
  \textbf{The second version has the scalars of two Higgs doublets.} Its
  two pair fields have eight real components, twice those of the Higgs
  doublet. In the pairing-only methods the response has four Goldstone
  modes and four massive modes, all four at \(\omega^{\rm H}_{\rm SH}\).
  Under the custodial symmetry each set is a singlet and a triplet.
  Besides the candidate for the Higgs boson there are therefore three
  more massive scalars, a neutral one and a charged one with its
  antiparticle, degenerate with it in these methods (Section 5.3).
\end{enumerate}

``I'' in this document is the assistant that does the work.

\hypertarget{the-hypotheses-and-the-ground-state}{%
\section{1. The hypotheses and the ground
state}\label{the-hypotheses-and-the-ground-state}}

\hypertarget{the-fields-and-their-charges}{%
\subsection{1.1 The fields and their
charges}\label{the-fields-and-their-charges}}

There are three spinor fields, all four-component Dirac fields of
\(SL(2,\mathbb C)\) and all singlets of colour:

\begin{itemize}
\tightlist
\item
  the isodoublet \(\Psi_{\rm LL}\) has eight components, numbered
  \(2d+s\), where \(d=0,\dots,3\) is the Dirac index of the chiral basis
  and \(s=0\) (up) or \(s=1\) (down) is the isospin index;
\item
  the singlets \(\chi_{\rm SH}\) and \(\chi'_{\rm SH}\) have four
  components each, numbered by the Dirac index.
\end{itemize}

The masses are
\(m_{\rm LL}=m^\chi_{\rm SH}=m^{\chi'}_{\rm SH}=0.2\,\Lambda_{\rm LL}\).
Together the three fields form one field of sixteen components,
\(\Psi_{\rm SH}=(\Psi_{\rm LL},\chi_{\rm SH},\chi'_{\rm SH})\):
components 0 to 7 are those of \(\Psi_{\rm LL}\), 8 to 11 those of
\(\chi_{\rm SH}\) and 12 to 15 those of \(\chi'_{\rm SH}\).

The gauge group is \(SL(2,\mathbb C)\times SU(2)_L\times U(1)_Y\). The
generators \(T^{a,{\rm SM}}\) (\(a=1,2,3\)) of \(SU(2)_L\) act as
\(\mathsf 1_4\otimes\tfrac12\tau^a\) on \(\Psi_{\rm LL}\) and as 0 on
the singlets, where \(\tau^a\) are the Pauli matrices. The hypercharge
\(Y^{\rm SM}\) is \(\tfrac14\) on \(\Psi_{\rm LL}\) and
\(\chi_{\rm SH}\) and \(-\tfrac34\) on \(\chi'_{\rm SH}\), and the
electric charge is \(Q^{\rm SM}=T^{3,{\rm SM}}+Y^{\rm SM}\). Two further
numbers are conserved, global symmetries that are not gauged:

\begin{itemize}
\tightlist
\item
  the particle number \(F^{\rm tot}_{\rm SH}\), 1 on every component;
\item
  the relative number \(R_{\rm SH}\), \(+1\) on the components of
  \(\Psi_{\rm LL}\) and \(-1\) on those of the singlets.
\end{itemize}

With equal masses there is one more symmetry, the custodial group
\(SU(2)_{\rm R}\). Its generators \(T^{a,{\rm R}}_{\rm SH}\) act as
\(\tfrac12\tau^a\) on the pair of singlets
\((\chi_{\rm SH},\chi'_{\rm SH})\), taken as a doublet with
\(\chi_{\rm SH}\) above, and as 0 on \(\Psi_{\rm LL}\). The hypercharge
is the combination
\(Y^{\rm SM}=T^{3,{\rm R}}_{\rm SH}+\tfrac14R_{\rm SH}\). The eight
generators \(T^{a,{\rm SM}}\), \(T^{a,{\rm R}}_{\rm SH}\),
\(R_{\rm SH}\) and \(F^{\rm tot}_{\rm SH}\) form the global group of the
hypotheses.

\sptkeptpenalty=10000
\begin{longtable}[]{@{}
  >{\raggedright\arraybackslash}p{(\columnwidth - 12\tabcolsep) * \real{0.1429}}
  >{\raggedright\arraybackslash}p{(\columnwidth - 12\tabcolsep) * \real{0.1429}}
  >{\raggedright\arraybackslash}p{(\columnwidth - 12\tabcolsep) * \real{0.1429}}
  >{\raggedright\arraybackslash}p{(\columnwidth - 12\tabcolsep) * \real{0.1429}}
  >{\raggedright\arraybackslash}p{(\columnwidth - 12\tabcolsep) * \real{0.1429}}
  >{\raggedright\arraybackslash}p{(\columnwidth - 12\tabcolsep) * \real{0.1429}}
  >{\raggedright\arraybackslash}p{(\columnwidth - 12\tabcolsep) * \real{0.1429}}@{}}
\toprule\noalign{}
\begin{minipage}[b]{\linewidth}\raggedright
field component
\end{minipage} & \begin{minipage}[b]{\linewidth}\raggedright
\(T^{3,{\rm SM}}\)
\end{minipage} & \begin{minipage}[b]{\linewidth}\raggedright
\(T^{3,{\rm R}}_{\rm SH}\)
\end{minipage} & \begin{minipage}[b]{\linewidth}\raggedright
\(Y^{\rm SM}\)
\end{minipage} & \begin{minipage}[b]{\linewidth}\raggedright
\(Q^{\rm SM}\)
\end{minipage} & \begin{minipage}[b]{\linewidth}\raggedright
\(F^{\rm tot}_{\rm SH}\)
\end{minipage} & \begin{minipage}[b]{\linewidth}\raggedright
\(R_{\rm SH}\)
\end{minipage} \\
\midrule\noalign{}
\endhead
\bottomrule\noalign{}
\endlastfoot
\(\Psi_{\rm LL}\), up & \(+\tfrac12\) & 0 & \(\tfrac14\) & \(\tfrac34\)
& 1 & \(+1\) \\*
\(\Psi_{\rm LL}\), down & \(-\tfrac12\) & 0 & \(\tfrac14\) &
\(-\tfrac14\) & 1 & \(+1\) \\*
\(\chi_{\rm SH}\) & 0 & \(+\tfrac12\) & \(\tfrac14\) & \(\tfrac14\) & 1
& \(-1\) \\*
\(\chi'_{\rm SH}\) & 0 & \(-\tfrac12\) & \(-\tfrac34\) & \(-\tfrac34\) &
1 & \(-1\) \\
\end{longtable}

Table 1.1. The charges of the particles of the three fields.
Antiparticles have the opposite charges.

\hypertarget{the-pair-fields}{%
\subsection{1.2 The pair fields}\label{the-pair-fields}}

In the chiral basis the Dirac matrices are
\(\gamma^0_{\rm LL}=\begin{pmatrix}0&\mathsf 1_2\\\mathsf 1_2&0\end{pmatrix}\),
\(\gamma^j_{\rm LL}=\begin{pmatrix}0&\sigma^{\rm LL}_j\\-\sigma^{\rm LL}_j&0\end{pmatrix}\)
(\(j=1,2,3\), with the Pauli matrices \(\sigma^{\rm LL}_j\)),
\(\beta_{\rm LL}=\gamma^0_{\rm LL}\),
\(\alpha^j_{\rm LL}=\gamma^0_{\rm LL}\gamma^j_{\rm LL}\),
\(\gamma^{\rm ch}_{\rm LL}=i\gamma^0_{\rm LL}\gamma^1_{\rm LL}\gamma^2_{\rm LL}\gamma^3_{\rm LL}\),
and the charge-conjugation matrix is
\(C_{\rm LL}=i\gamma^2_{\rm LL}\gamma^0_{\rm LL}\). The pair operators
of \(\Psi_{\rm LL}\) with the singlets are
\[\Pi^s_{\rm SH}=\tfrac12\Psi^T_{\rm LL}\big(C_{\rm LL}\gamma^{\rm ch}_{\rm LL}\otimes\mathbf e^s_{\rm SH}\big)\chi_{\rm SH},\qquad
\Pi'^s_{\rm SH}=\tfrac12\Psi^T_{\rm LL}\big(C_{\rm LL}\gamma^{\rm ch}_{\rm LL}\otimes\mathbf e^s_{\rm SH}\big)\chi'_{\rm SH},\qquad s={\rm u},{\rm d},\tag{1.1}\]
with the isospin columns \(\mathbf e^{\rm u}_{\rm SH}=(1,0)^T\) and
\(\mathbf e^{\rm d}_{\rm SH}=(0,1)^T\). Each annihilates two particles
in a state of spin zero and positive parity. On the whole field,
\(\Psi^T_{\rm SH}\mathsf K^s_{\rm SH}\Psi_{\rm SH}=2\Pi^s_{\rm SH}\) and
\(\Psi^T_{\rm SH}\mathsf K'^s_{\rm SH}\Psi_{\rm SH}=2\Pi'^s_{\rm SH}\).
Here \(\mathsf K^s_{\rm SH}\) is the antisymmetric \(16\times16\) matrix
with the block
\(\tfrac12C_{\rm LL}\gamma^{\rm ch}_{\rm LL}\otimes\mathbf e^s_{\rm SH}\)
in the rows of \(\Psi_{\rm LL}\) and the columns of \(\chi_{\rm SH}\),
minus its transpose in the opposite block, and zero elsewhere;
\(\mathsf K'^s_{\rm SH}\) is the same with the columns of
\(\chi'_{\rm SH}\). These four matrices are real.

The pair fields are the expectation values:
\(\Phi_{\rm SH}=(\langle\Pi^{\rm u}_{\rm SH}\rangle,\langle\Pi^{\rm d}_{\rm SH}\rangle)\)
is a doublet of \(SU(2)_L\) with the hypercharge \(\tfrac12\), the
quantum numbers of the Higgs doublet, and
\(\Phi'_{\rm SH}=(\langle\Pi'^{\rm u}_{\rm SH}\rangle,\langle\Pi'^{\rm d}_{\rm SH}\rangle)\)
is a doublet with the hypercharge \(-\tfrac12\). Their neutral
components are \(\langle\Pi^{\rm d}_{\rm SH}\rangle\) and
\(\langle\Pi'^{\rm u}_{\rm SH}\rangle\).

\hypertarget{the-hamiltonian-density-and-the-methods}{%
\subsection{1.3 The Hamiltonian density and the
methods}\label{the-hamiltonian-density-and-the-methods}}

The Hamiltonian density is
\[\begin{aligned}\mathcal H_{\rm SH}={}&\Psi^\dagger_{\rm SH}\mathsf h_{\rm SH}(-i\nabla)\Psi_{\rm SH}+\tfrac12G^{\rm S}_{\rm SF}{:}(\bar\Psi_{\rm LL}\Psi_{\rm LL})^2{:}
+\tfrac12G^{\rm V}_{\rm SF}{:}\Big[(\Psi^\dagger_{\rm LL}\Psi_{\rm LL})^2-\sum_j\big(\Psi^\dagger_{\rm LL}(\alpha^j_{\rm LL}\otimes\mathsf 1_2)\Psi_{\rm LL}\big)^2\Big]{:}\\
&-G_{\rm SH}\sum_{s={\rm u},{\rm d}}\big[{:}(\Pi^s_{\rm SH})^\dagger\Pi^s_{\rm SH}{:}+{:}(\Pi'^s_{\rm SH})^\dagger\Pi'^s_{\rm SH}{:}\big].\end{aligned}\tag{1.2}\]
Its parts:

\begin{itemize}
\tightlist
\item
  \(\mathsf h_{\rm SH}(\vec q_{\rm LL})\) is the free Dirac Hamiltonian
  of the sixteen components on a plane wave of momentum
  \(\vec q_{\rm LL}\), block-diagonal with
  \((\alpha^j_{\rm LL}q^j_{\rm LL}+\beta_{\rm LL}m_{\rm LL})\otimes\mathsf 1_2\)
  on \(\Psi_{\rm LL}\) and
  \(\alpha^j_{\rm LL}q^j_{\rm LL}+\beta_{\rm LL}m_{\rm LL}\) on each
  singlet (summed over \(j\)); the momenta are those with
  \(\lvert\vec q_{\rm LL}\rvert<\Lambda_{\rm LL}\);
\item
  \(\bar\Psi_{\rm LL}=\Psi^\dagger_{\rm LL}(\gamma^0_{\rm LL}\otimes\mathsf 1_2)\),
  and the colons order the products of field operators so that their
  expectation value in the normal vacuum (the state without pairs) is
  zero;
\item
  the couplings are those of the potential \(V^{(1)}_{\rm SH}\) of stage
  3: \(G^{\rm S}_{\rm SF}=-8\) and \(G^{\rm V}_{\rm SF}=0.8\) (the
  scalar and the vector contact terms of the author's ansatz) and
  \(G_{\rm SH}=100.8\) (the pairing terms), in units of
  \(\Lambda_{\rm LL}^{-2}\).
\end{itemize}

In the \textbf{first version} of the hypotheses the terms of
\(\Pi'^s_{\rm SH}\) are absent, and \(\chi'_{\rm SH}\) is a free field.

The three methods approximate the state by a uniform quasi-free state.
Such a state is the ground state of a quadratic Hamiltonian with
constant self-energies, a Hermitian \(16\times16\) matrix
\(\mathsf N_{\rm SH}\) and an antisymmetric one \(\mathsf P_{\rm SH}\).
At the momentum \(\vec q_{\rm LL}\) its Bogoliubov-de Gennes (BdG)
matrix is
\[\mathsf H_{\rm SH}(\vec q_{\rm LL})=\begin{pmatrix}\mathsf h_{\rm SH}(\vec q_{\rm LL})+\mathsf N_{\rm SH}&\mathsf P_{\rm SH}\\\mathsf P_{\rm SH}^\dagger&-\big(\mathsf h_{\rm SH}(-\vec q_{\rm LL})+\mathsf N_{\rm SH}\big)^T\end{pmatrix},\tag{1.3}\]
and the state fills its negative eigenvalues. Its local matrices are the
integrals over the momenta of the normal correlation matrix, measured
from the normal vacuum, and of the pair correlation matrix:
\(\varrho_{\rm SH}\) and \(\kappa_{\rm SH}\). The energy per volume of
the state, measured from the normal vacuum, is
\(\mathcal E^X_{\rm SH}\): the kinetic part plus the interaction energy
that method \(X\) assigns to the local matrices. A method keeps the
three Wick contractions of every contact term with its own weights:

\begin{itemize}
\tightlist
\item
  HS rewrites each four-fermion term as the exchange of an auxiliary
  field and keeps the direct contraction of the terms decoupled in
  density channels and the pair contraction of those decoupled in pair
  channels. Its parameters \(s^{\rm HS,S}_{\rm SF}\) and
  \(s^{\rm HS,V}_{\rm SF}\) are the shares of the scalar and of the
  vector term decoupled in density channels; HS (1, 1), HS (0, 1) and HS
  (0, 0) denote \((s^{\rm HS,S}_{\rm SF},s^{\rm HS,V}_{\rm SF})\). The
  pairing terms are decoupled in their own pair channels. HS keeps no
  exchange contraction.
\item
  KS keeps all three contractions, the exchange evaluated with the
  projection of \(\varrho_{\rm SH}\) onto sixteen local densities: the
  number and the scalar densities of \(\Psi_{\rm LL}\) as \(2\times2\)
  matrices in the isospin space and those of the pair of singlets as
  \(2\times2\) matrices in the space of the two singlets.
\item
  CJT keeps all three contractions with the weight 1 (the two-loop
  effective action of the propagator).
\end{itemize}

The self-energies are the derivatives of the interaction energy with
respect to the local matrices, and a stationary state is a state whose
self-energies are those that its own local matrices produce. The
programs find stationary states by iteration on a grid of 24 radial
nodes of \(\lvert\vec q_{\rm LL}\rvert\) times the six directions of the
coordinate axes.

\hypertarget{the-ground-states-of-stage-3}{%
\subsection{1.4 The ground states of stage
3}\label{the-ground-states-of-stage-3}}

In every method the ground state of \(V^{(1)}_{\rm SH}\) in the second
version is the custodial Higgs superfluid. Its pair self-energy is
\(\mathsf P_{\rm SH}=d_{\rm SH}\,\mathsf P^{\rm vac}_{\rm SH}\) with a
real number \(d_{\rm SH}=d^0_{\rm SH}>0\) and the direction of the
vacuum
\[\mathsf P^{\rm vac}_{\rm SH}=(\mathsf K'^{\rm u}_{\rm SH})^\dagger-(\mathsf K^{\rm d}_{\rm SH})^\dagger\qquad\text{(second version)},\qquad
\mathsf P^{\rm vac}_{\rm SH}=(\mathsf K^{\rm d}_{\rm SH})^\dagger\qquad\text{(first version)}.\tag{1.4}\]
A rotation by the particle number multiplies \(\mathsf P_{\rm SH}\) by a
phase. The programs rotate that phase away, so that \(d_{\rm SH}\) is
real and positive. In the second version the down particles are paired
with \(\chi_{\rm SH}\) and the up particles with \(\chi'_{\rm SH}\), the
two pair fields have equal lengths, and the state keeps the electric
charge, the custodial group \(SU(2)_{\rm V}\) (generated by
\(T^{a,{\rm SM}}+T^{a,{\rm R}}_{\rm SH}\)) and \(R_{\rm SH}\). In the
first version the ground state is the condensate of \(\Phi_{\rm SH}\)
alone: the down particles are paired with \(\chi_{\rm SH}\) and the up
particles are unpaired.

The \textbf{pair amplitude} of a state is
\[\varphi_{\rm SH}=\lvert\langle\Pi^{\rm d}_{\rm SH}\rangle\rvert=\tfrac12\big\lvert\mathrm{tr}\big(\mathsf K^{\rm d}_{\rm SH}\kappa_{\rm SH}^T\big)\big\rvert,\tag{1.5}\]
the length of the neutral component of \(\Phi_{\rm SH}\). In the states
of the second version computed here it equals
\(\lvert\langle\Pi'^{\rm u}_{\rm SH}\rangle\rvert\), the length of the
neutral component of \(\Phi'_{\rm SH}\) (check H-2d); the program uses
the mean of the two.

\hypertarget{the-effective-potential}{%
\section{2. The effective potential}\label{the-effective-potential}}

\hypertarget{definition-and-method}{%
\subsection{2.1 Definition and method}\label{definition-and-method}}

The effective potential \(V^{\rm eff}_{\rm SH}(\varphi_{\rm SH})\) is
the lowest energy per volume of the uniform states with the pair
amplitude \(\varphi_{\rm SH}\) along the direction of the vacuum. It is
computed as follows. Minimizing the energy at a fixed pair amplitude
adds a Lagrange multiplier times
\(\Pi^{\rm d}_{\rm SH}+(\Pi^{\rm d}_{\rm SH})^\dagger\) (and the same
for \(\Pi'^{\rm u}_{\rm SH}\) in the second version) to the Hamiltonian.
The stationary state of the constrained problem is therefore the ground
state of (1.3) with a pair self-energy along
\(\mathsf P^{\rm vac}_{\rm SH}\),
\(\mathsf P_{\rm SH}=d_{\rm SH}\mathsf P^{\rm vac}_{\rm SH}\), for some
\(d_{\rm SH}\). The normal self-energy \(\mathsf N_{\rm SH}\) is the one
that its own local matrices produce. So the program:

\begin{itemize}
\tightlist
\item
  fixes \(d_{\rm SH}\) and iterates \(\mathsf N_{\rm SH}\) to
  self-consistency;
\item
  records the energy \(\mathcal E^X_{\rm SH}\) and the pair amplitude
  \(\varphi_{\rm SH}\) of the state;
\item
  does this for 33 values
  \(d_{\rm SH}=0,0.05\,d^0_{\rm SH},\dots,1.6\,d^0_{\rm SH}\).
\end{itemize}

The curve
\((\varphi_{\rm SH}(d_{\rm SH}),\mathcal E^X_{\rm SH}(d_{\rm SH}))\) is
\(V^{\rm eff}_{\rm SH}\) in parametric form. At
\(d_{\rm SH}=d^0_{\rm SH}\) the pair self-energy that the local matrices
produce equals the fixed one: the state is the ground state (check H-1).
The derivatives at the minimum
\(\varphi^0_{\rm SH}=\varphi_{\rm SH}(d^0_{\rm SH})\) follow from
central differences in \(d_{\rm SH}\) with the steps 0.01 and 0.02,
combined by Richardson extrapolation:
\[V^{\rm eff\,\prime}_{\rm SH}=\frac{\mathcal E'}{\varphi'},\qquad V^{\rm eff\,\prime\prime}_{\rm SH}=\frac{\mathcal E''-V^{\rm eff\,\prime}_{\rm SH}\varphi''}{\varphi'^2},\tag{2.1}\]
where the primes on \(\mathcal E\) and \(\varphi\) are derivatives with
respect to \(d_{\rm SH}\).

\hypertarget{the-closed-form-of-the-pairing-only-methods}{%
\subsection{2.2 The closed form of the pairing-only
methods}\label{the-closed-form-of-the-pairing-only-methods}}

In HS (0, 1) and HS (0, 0) only the pairing terms act on these states,
and the normal self-energy is zero. Then every quasi-particle has the
energy \(E=(q_{\rm LL}^2+(M^{\rm eff}_{\rm SH})^2)^{1/2}\) with
\(M^{\rm eff}_{\rm SH}=(m_{\rm LL}^2+d_{\rm SH}^2/4)^{1/2}\), by the
identity
\((C_{\rm LL}\gamma^{\rm ch}_{\rm LL})(-\alpha^j_{\rm LL}q^j_{\rm LL}+\beta_{\rm LL}m_{\rm LL})(C_{\rm LL}\gamma^{\rm ch}_{\rm LL})^{-1}=(\alpha^j_{\rm LL}q^j_{\rm LL}+\beta_{\rm LL}m_{\rm LL})^T\),
and the pair correlation matrix is \(\mathsf P_{\rm SH}/(2E)\). With
\[I^{(1)}_{\rm SF}(M)=\int_0^{\Lambda_{\rm LL}}\frac{q^2\,dq}{(q^2+M^2)^{1/2}},\qquad J^{(3)}_{\rm SH}(M)=\int_0^{\Lambda_{\rm LL}}\frac{q^2\,dq}{(q^2+M^2)^{3/2}}=\mathrm{arsinh}\frac{\Lambda_{\rm LL}}{M}-\frac{\Lambda_{\rm LL}}{(\Lambda_{\rm LL}^2+M^2)^{1/2}}\tag{2.2}\]
this gives, for each condensed pair,
\[\varphi_{\rm SH}(d_{\rm SH})=\frac{d_{\rm SH}\,I^{(1)}_{\rm SF}(M^{\rm eff}_{\rm SH})}{4\pi^2},\qquad\frac{dV^{\rm eff}_{\rm SH}}{d\varphi_{\rm SH}}=2\big(d_{\rm SH}-G_{\rm SH}\varphi_{\rm SH}\big)\ \text{(one pair)}.\tag{2.3}\]
The second relation holds because the BdG ground state with the pair
self-energy \(d_{\rm SH}\mathsf P^{\rm vac}_{\rm SH}\) minimizes the
kinetic energy minus \(2d_{\rm SH}\varphi_{\rm SH}\), and the
interaction energy of one pair is \(-G_{\rm SH}\varphi_{\rm SH}^2\). The
minimum is at \(d_{\rm SH}=G_{\rm SH}\varphi_{\rm SH}\), which is the
gap equation
\(G_{\rm SH}I^{(1)}_{\rm SF}(M^{\rm eff}_{\rm SH})/(4\pi^2)=1\), and
there \(\varphi^0_{\rm SH}=d^0_{\rm SH}/G_{\rm SH}\). Differentiating
once more, with \(dI^{(1)}_{\rm SF}/dM=-M\,J^{(3)}_{\rm SH}(M)\),
\[V^{\rm eff\,\prime\prime}_{\rm SH}(\varphi^0_{\rm SH})=\frac{2\pi^2(d^0_{\rm SH})^2J^{(3)}_{\rm SH}}{I^{(1)}_{\rm SF}\big(I^{(1)}_{\rm SF}-(d^0_{\rm SH})^2J^{(3)}_{\rm SH}/4\big)}\tag{2.4}\]
for one pair, at \(M=M^{\rm eff}_{\rm SH}\). In the second version both
pairs move together along \(\mathsf P^{\rm vac}_{\rm SH}\), and the
curvature is twice (2.4). The Wolfram script evaluates (2.4) to 30
digits (check WH-3) and the Python program computes it from the energies
of its own engine (check YH-3); both agree with the Rust program.

\hypertarget{the-results}{%
\subsection{2.3 The results}\label{the-results}}

Table 2.1 gives, for each method and version, the minimum of the
effective potential: the pair self-energy and the pair amplitude of the
ground state, its energy, and the curvature there.

Table 2.1. The minimum of the effective potential of
\(V^{(1)}_{\rm SH}\).

\sptkeptpenalty=10000
\begin{longtable}[]{@{}
  >{\raggedright\arraybackslash}p{(\columnwidth - 10\tabcolsep) * \real{0.1176}}
  >{\raggedright\arraybackslash}p{(\columnwidth - 10\tabcolsep) * \real{0.1765}}
  >{\raggedright\arraybackslash}p{(\columnwidth - 10\tabcolsep) * \real{0.1765}}
  >{\raggedright\arraybackslash}p{(\columnwidth - 10\tabcolsep) * \real{0.1765}}
  >{\raggedright\arraybackslash}p{(\columnwidth - 10\tabcolsep) * \real{0.1765}}
  >{\raggedright\arraybackslash}p{(\columnwidth - 10\tabcolsep) * \real{0.1765}}@{}}
\toprule\noalign{}
\begin{minipage}[b]{\linewidth}\raggedright
version
\end{minipage} & \begin{minipage}[b]{\linewidth}\raggedright
method
\end{minipage} & \begin{minipage}[b]{\linewidth}\raggedright
\(d^0_{\rm SH}\)
\end{minipage} & \begin{minipage}[b]{\linewidth}\raggedright
\(\varphi^0_{\rm SH}\)
\end{minipage} & \begin{minipage}[b]{\linewidth}\raggedright
\(10^4\,\mathcal E^X_{\rm SH}\)
\end{minipage} & \begin{minipage}[b]{\linewidth}\raggedright
\(V^{\rm eff\,\prime\prime}_{\rm SH}(\varphi^0_{\rm SH})\)
\end{minipage} \\
\midrule\noalign{}
\endhead
\bottomrule\noalign{}
\endlastfoot
second & HS (1, 1) & \(0.832878\) & \(0.00826267\) & \(-11.2952\) &
\(149.188\) \\*
second & HS (0, 1) & \(0.812846\) & \(0.00806394\) & \(-10.5678\) &
\(143.361\) \\*
second & HS (0, 0) & \(0.812846\) & \(0.00806394\) & \(-10.5678\) &
\(143.361\) \\*
second & KS & \(0.890787\) & \(0.00883717\) & \(-16.4192\) &
\(181.683\) \\*
second & CJT & \(0.890787\) & \(0.00883717\) & \(-16.4192\) &
\(181.683\) \\*
first & HS (1, 1) & \(0.82581\) & \(0.00819256\) & \(-5.49853\) &
\(73.148\) \\*
first & HS (0, 1) & \(0.812846\) & \(0.00806394\) & \(-5.28389\) &
\(71.6803\) \\*
first & HS (0, 0) & \(0.812846\) & \(0.00806394\) & \(-5.28389\) &
\(71.6803\) \\*
first & KS & \(0.867286\) & \(0.00860402\) & \(-6.2857\) &
\(80.0121\) \\*
first & CJT & \(0.867286\) & \(0.00860402\) & \(-6.2857\) &
\(80.0121\) \\
\end{longtable}

In every method and both versions the curve is a Mexican hat (check
H-3): \(V^{\rm eff}_{\rm SH}=0\) at \(\varphi_{\rm SH}=0\), it decreases
to its minimum at \(d_{\rm SH}=d^0_{\rm SH}\) and increases beyond it,
and \(\varphi_{\rm SH}\) grows monotonically with \(d_{\rm SH}\). The
curvature is positive and the first derivative vanishes at
\(\varphi^0_{\rm SH}\) to the accuracy of the differences (checks H-2a
and H-2b). In HS (0, 1) and HS (0, 0) the curvature of the second
version is twice that of the first, because two equal pairs move
together (Section 2.2); in the other methods the ratio is \(2.04\) to
\(2.27\).

\hypertarget{the-susceptibilities-and-the-masses-of-the-gauge-bosons}{%
\section{3. The susceptibilities and the masses of the gauge
bosons}\label{the-susceptibilities-and-the-masses-of-the-gauge-bosons}}

\hypertarget{chemical-potentials-as-the-time-components-of-the-gauge-fields}{%
\subsection{3.1 Chemical potentials as the time components of the gauge
fields}\label{chemical-potentials-as-the-time-components-of-the-gauge-fields}}

A uniform, static gauge field couples to the charge densities. Its time
component \(A^a_0\) acts on the state like a chemical potential
\(\mu^a_{\rm SH}=g_aA^a_0\) of the charge \(Q_a\), with \(g_a\) the
coupling of the generator. In the effective Lagrangian of the pair
fields the term \(Z_{\rm SH}\lvert D_\mu\Phi\rvert^2\) (with \(D_\mu\)
the covariant derivative and \(Z_{\rm SH}\) the kinetic coefficient)
gives the gauge fields the mass term
\(\tfrac12\sum_{ab}g_ag_b\chi^{ab}_{\rm SH}A^a_\mu A^{b,\mu}\). Here
\(\chi^{ab}_{\rm SH}\) is the susceptibility of the charges to the
chemical potentials:
\[\chi^{ab}_{\rm SH}=\frac{\partial\langle Q_a\rangle}{\partial\mu^b_{\rm SH}}\Big\rvert_{\mu=0}=-\frac{\partial^2\mathcal E^X_{\rm SH}}{\partial\mu^a_{\rm SH}\partial\mu^b_{\rm SH}}\Big\rvert_{\mu=0},\tag{3.1}\]
where \(\mathcal E^X_{\rm SH}(\mu)\) is the energy of the ground state
of
\(\mathcal H_{\rm SH}-\sum_b\mu^b_{\rm SH}\Psi^\dagger_{\rm SH}Q_b\Psi_{\rm SH}\)
and \(\langle Q_a\rangle=\mathrm{tr}(Q_a\varrho_{\rm SH})\) its charge
density. The two forms are equal by the theorem of Hellmann and Feynman:
at a stationary state the derivative of the energy with respect to
\(\mu^a_{\rm SH}\) is \(-\langle Q_a\rangle\). For a Lorentz-invariant
vacuum the mass of the time component equals that of the space
components. The relation of \(\chi^{ab}_{\rm SH}\) to the masses uses
this.

\textbf{Why the time components.} The normal vacuum has a gap in its
spectrum and keeps every charge, so a small chemical potential changes
nothing in it: its susceptibility is exactly zero, and the gauge bosons
get no mass without a condensate (check YH-9: the charge densities of
the normal vacuum stay zero to \(3\times10^{-15}\) of the chemical
potential). The response to the space components, the currents, is not
zero in the normal vacuum with the sharp cutoff of the momenta. A
uniform vector potential shifts the momenta, and the current of the
shifted Dirac sea is a surface term of the cutoff: an artifact, since a
uniform vector potential is a gauge transformation. The time components
avoid it.

\textbf{The computation.} For each of the eight generators \(Q_b\),
which are \(T^{1,{\rm SM}}\), \(T^{2,{\rm SM}}\), \(T^{3,{\rm SM}}\),
\(T^{1,{\rm R}}_{\rm SH}\), \(T^{2,{\rm R}}_{\rm SH}\),
\(T^{3,{\rm R}}_{\rm SH}\), \(R_{\rm SH}\) and \(F^{\rm tot}_{\rm SH}\),
the program solves for the ground state of
\(\mathcal H_{\rm SH}-\mu^b_{\rm SH}Q_b\) at \(\mu^b_{\rm SH}=\pm0.005\)
and \(\pm0.01\). It starts from the ground state and iterates both
self-energies to self-consistency. It records the eight charge densities
and forms central differences, combined by Richardson extrapolation.
This gives the \(8\times8\) matrix \(\chi^{ab}_{\rm SH}\). By (3.1) it
must be symmetric (check H-4a). The rows of the hypercharge and of the
electric charge follow by linearity from
\(Y^{\rm SM}=T^{3,{\rm R}}_{\rm SH}+\tfrac14R_{\rm SH}\) and
\(Q^{\rm SM}=T^{3,{\rm SM}}+Y^{\rm SM}\). The Python program computes
the same susceptibilities by the second form of (3.1), from the energies
of its own engine (check YH-2).

\hypertarget{the-structure-of-the-susceptibilities}{%
\subsection{3.2 The structure of the
susceptibilities}\label{the-structure-of-the-susceptibilities}}

Write \(\chi^{11}_{\rm SH}\), \(\chi^{33}_{\rm SH}\) for the
susceptibilities of \(T^{1,{\rm SM}}\) and \(T^{3,{\rm SM}}\),
\(\chi^{YY}_{\rm SH}\) for that of \(Y^{\rm SM}\),
\(\chi^{3Y}_{\rm SH}\) for the mixed one of \(T^{3,{\rm SM}}\) and
\(Y^{\rm SM}\), and \(\chi^{FF}_{\rm SH}\) for that of the particle
number. The computed matrices have the following structure, in every
method:

\begin{itemize}
\tightlist
\item
  \textbf{The unbroken generators do not respond.} In the second version
  the generators that stage 3 found unbroken,
  \(T^{a,{\rm SM}}+T^{a,{\rm R}}_{\rm SH}\) and \(R_{\rm SH}\), have
  zero susceptibility with every generator (check H-4e). The electric
  charge does not respond in either version (check H-4g).
\item
  \textbf{The custodial structure.} In the second version
  \(\chi(T^{a,{\rm SM}},T^{a,{\rm SM}})=\chi(T^{a,{\rm R}}_{\rm SH},T^{a,{\rm R}}_{\rm SH})=-\chi(T^{a,{\rm SM}},T^{a,{\rm R}}_{\rm SH})\),
  equal for \(a=1,2,3\) (check H-4f). This is the structure that the
  effective Lagrangian
  \(Z_{\rm SH}\,\mathrm{tr}((\partial\mathcal M_{\rm SH})^\dagger\partial\mathcal M_{\rm SH})\)
  of the bidoublet
  \(\mathcal M_{\rm SH}=(\Phi_{\rm SH},\Phi'_{\rm SH})\) (rows: the
  isospin components; columns: the two singlets) gives at
  \(\mathcal M_{\rm SH}\propto i\tau^2\) (check WH-1, symbolic). It
  implies
  \(\chi^{11}_{\rm SH}=\chi^{33}_{\rm SH}=\chi^{YY}_{\rm SH}=-\chi^{3Y}_{\rm SH}\).
\item
  \textbf{The number of broken generators.} The symmetric matrix
  \(\chi^{ab}_{\rm SH}\) has exactly four nonzero eigenvalues in the
  second version, all positive. In the first version, among the
  generators that are symmetries there, \(T^{1,{\rm SM}}\),
  \(T^{2,{\rm SM}}\), \(T^{3,{\rm SM}}\), \(R_{\rm SH}\) and
  \(F^{\rm tot}_{\rm SH}\), it has three (check H-4d). These are the
  numbers of broken generators.
\end{itemize}

Table 3.1. The susceptibilities of the ground states of
\(V^{(1)}_{\rm SH}\) (times \(10^3\)), and
\(\rho_{\rm SH}=\chi^{11}_{\rm SH}/\chi^{33}_{\rm SH}\).

\sptkeptpenalty=10000
\begin{longtable}[]{@{}
  >{\raggedright\arraybackslash}p{(\columnwidth - 14\tabcolsep) * \real{0.0976}}
  >{\raggedright\arraybackslash}p{(\columnwidth - 14\tabcolsep) * \real{0.1463}}
  >{\raggedright\arraybackslash}p{(\columnwidth - 14\tabcolsep) * \real{0.1220}}
  >{\raggedright\arraybackslash}p{(\columnwidth - 14\tabcolsep) * \real{0.1220}}
  >{\raggedright\arraybackslash}p{(\columnwidth - 14\tabcolsep) * \real{0.1220}}
  >{\raggedright\arraybackslash}p{(\columnwidth - 14\tabcolsep) * \real{0.1220}}
  >{\raggedright\arraybackslash}p{(\columnwidth - 14\tabcolsep) * \real{0.1220}}
  >{\raggedright\arraybackslash}p{(\columnwidth - 14\tabcolsep) * \real{0.1463}}@{}}
\toprule\noalign{}
\begin{minipage}[b]{\linewidth}\raggedright
version
\end{minipage} & \begin{minipage}[b]{\linewidth}\raggedright
method
\end{minipage} & \begin{minipage}[b]{\linewidth}\raggedright
\(\chi^{11}_{\rm SH}\)
\end{minipage} & \begin{minipage}[b]{\linewidth}\raggedright
\(\chi^{33}_{\rm SH}\)
\end{minipage} & \begin{minipage}[b]{\linewidth}\raggedright
\(\chi^{YY}_{\rm SH}\)
\end{minipage} & \begin{minipage}[b]{\linewidth}\raggedright
\(\chi^{3Y}_{\rm SH}\)
\end{minipage} & \begin{minipage}[b]{\linewidth}\raggedright
\(\chi^{FF}_{\rm SH}\)
\end{minipage} & \begin{minipage}[b]{\linewidth}\raggedright
\(\rho_{\rm SH}\)
\end{minipage} \\
\midrule\noalign{}
\endhead
\bottomrule\noalign{}
\endlastfoot
second & HS (1, 1) & \(2.73681\) & \(2.73681\) & \(2.73681\) &
\(-2.73681\) & \(43.4088\) & \(1\) \\*
second & HS (0, 1) & \(2.60214\) & \(2.60214\) & \(2.60214\) &
\(-2.60214\) & \(41.2905\) & \(1\) \\*
second & HS (0, 0) & \(2.60214\) & \(2.60214\) & \(2.60214\) &
\(-2.60214\) & \(41.6343\) & \(1\) \\*
second & KS & \(3.11209\) & \(3.11209\) & \(3.11209\) & \(-3.11209\) &
\(58.367\) & \(1\) \\*
second & CJT & \(3.11209\) & \(3.11209\) & \(3.11209\) & \(-3.11209\) &
\(58.367\) & \(1\) \\*
first & HS (1, 1) & \(1.66832\) & \(1.33798\) & \(1.33798\) &
\(-1.33798\) & \(21.4077\) & \(1.2468945\) \\*
first & HS (0, 1) & \(1.59944\) & \(1.29568\) & \(1.29568\) &
\(-1.29568\) & \(20.7308\) & \(1.2344421\) \\*
first & HS (0, 0) & \(1.59944\) & \(1.30107\) & \(1.30107\) &
\(-1.30107\) & \(20.8171\) & \(1.2293239\) \\*
first & KS & \(1.88651\) & \(1.56984\) & \(1.56984\) & \(-1.56984\) &
\(25.1174\) & \(1.2017215\) \\*
first & CJT & \(1.88651\) & \(1.56984\) & \(1.56984\) & \(-1.56984\) &
\(25.1174\) & \(1.2017215\) \\
\end{longtable}

\hypertarget{the-masses-of-the-gauge-bosons-and-the-parameter-rho}{%
\subsection{\texorpdfstring{3.3 The masses of the gauge bosons and the
parameter
\(\rho\)}{3.3 The masses of the gauge bosons and the parameter \textbackslash rho}}\label{the-masses-of-the-gauge-bosons-and-the-parameter-rho}}

The mass matrix of the four electroweak gauge fields
\((W^1_\mu,W^2_\mu,W^3_\mu,B_\mu)\) is
\[\big(\mathsf M^{\rm g}_{\rm SH}\big)_{ab}=g_ag_b\,\chi^{ab}_{\rm SH},\qquad(g_a)=(g,g,g,g'),\tag{3.2}\]
with \(a,b\) running over \(T^{1,{\rm SM}}\), \(T^{2,{\rm SM}}\),
\(T^{3,{\rm SM}}\) and \(Y^{\rm SM}\). The couplings are those of the
measured inputs of stage 2: \(g=0.652771\) and \(g'=0.350046\). Its
eigenvalues are, in every method and version (checks H-5a, H-5b):

\begin{itemize}
\tightlist
\item
  zero, with the eigenvector \((0,0,g',g)/(g^2+g'^2)^{1/2}\): the photon
  is massless;
\item
  \(M_{W,{\rm SH}}^2=g^2\chi^{11}_{\rm SH}\), twice: \(W^\pm_\mu\);
\item
  \(M_{Z,{\rm SH}}^2\): \(Z_\mu\).
\end{itemize}

The scale of the breaking is defined as in the Standard Model by
\(M_{W,{\rm SH}}=g\,v_{\rm SH}/2\), so
\[v_{\rm SH}=2\big(\chi^{11}_{\rm SH}\big)^{1/2},\qquad\rho_{\rm SH}=\frac{M_{W,{\rm SH}}^2}{M_{Z,{\rm SH}}^2\cos^2\theta_W}=\frac{\chi^{11}_{\rm SH}}{\chi^{33}_{\rm SH}},\tag{3.3}\]
with \(\cos^2\theta_W=g^2/(g^2+g'^2)\) (check H-6a). The second form of
\(\rho_{\rm SH}\) follows from the structure of Section 3.2: the zero
row of \(Q^{\rm SM}\) gives \(\chi^{3Y}_{\rm SH}=-\chi^{33}_{\rm SH}\)
and \(\chi^{YY}_{\rm SH}=\chi^{33}_{\rm SH}\).

In the second version \(\rho_{\rm SH}\) is \(1\) in every method (check
H-6b): the custodial symmetry of the vacuum holds to the accuracy of the
computation. The measured value, with the weak mixing angle of the
\(\overline{\rm MS}\) scheme, is \(\rho=1.0103\). Its deviation from 1
comes from the radiative corrections of the Standard Model, which the
mean field leaves out. In the first version \(\rho_{\rm SH}\) is
\(1.202\) to \(1.247\). There the condensate pairs only the down
component, and the rotations of \(T^{1,{\rm SM}}\) and
\(T^{2,{\rm SM}}\), which turn a paired down particle into an unpaired
up one, cost a different energy than the phase rotation of
\(T^{3,{\rm SM}}\). That is the reason for the second singlet.

\hypertarget{the-scale-in-gev}{%
\section{4. The scale in GeV}\label{the-scale-in-gev}}

Setting \(v_{\rm SH}=v=246.22\) GeV, the value of the Standard Model
from the Fermi constant, fixes the cutoff in GeV:
\(\Lambda_{\rm LL}=v/v_{\rm SH}\) when \(v_{\rm SH}\) is in units of the
cutoff. Then \(M_{W,{\rm SH}}=gv/2=80.3625\) GeV in every method, by
construction, and
\(M_{Z,{\rm SH}}=M_{W,{\rm SH}}/(\rho_{\rm SH}^{1/2}\cos\theta_W)\).

Table 4.1. The scale and the masses of the gauge bosons in GeV.

\sptkeptpenalty=10000
\begin{longtable}[]{@{}
  >{\raggedright\arraybackslash}p{(\columnwidth - 12\tabcolsep) * \real{0.1000}}
  >{\raggedright\arraybackslash}p{(\columnwidth - 12\tabcolsep) * \real{0.1500}}
  >{\raggedright\arraybackslash}p{(\columnwidth - 12\tabcolsep) * \real{0.1500}}
  >{\raggedright\arraybackslash}p{(\columnwidth - 12\tabcolsep) * \real{0.1500}}
  >{\raggedright\arraybackslash}p{(\columnwidth - 12\tabcolsep) * \real{0.1500}}
  >{\raggedright\arraybackslash}p{(\columnwidth - 12\tabcolsep) * \real{0.1500}}
  >{\raggedright\arraybackslash}p{(\columnwidth - 12\tabcolsep) * \real{0.1500}}@{}}
\toprule\noalign{}
\begin{minipage}[b]{\linewidth}\raggedright
version
\end{minipage} & \begin{minipage}[b]{\linewidth}\raggedright
method
\end{minipage} & \begin{minipage}[b]{\linewidth}\raggedright
\(v_{\rm SH}\) (units of \(\Lambda_{\rm LL}\))
\end{minipage} & \begin{minipage}[b]{\linewidth}\raggedright
\(\Lambda_{\rm LL}\) (GeV)
\end{minipage} & \begin{minipage}[b]{\linewidth}\raggedright
\(M_{W,{\rm SH}}\) (GeV)
\end{minipage} & \begin{minipage}[b]{\linewidth}\raggedright
\(M_{Z,{\rm SH}}\) (GeV)
\end{minipage} & \begin{minipage}[b]{\linewidth}\raggedright
\(\rho_{\rm SH}\)
\end{minipage} \\
\midrule\noalign{}
\endhead
\bottomrule\noalign{}
\endlastfoot
second & HS (1, 1) & \(0.104629\) & \(2353.3\) & \(80.3625\) &
\(91.1879\) & \(1\) \\*
second & HS (0, 1) & \(0.102022\) & \(2413.4\) & \(80.3625\) &
\(91.1879\) & \(1\) \\*
second & HS (0, 0) & \(0.102022\) & \(2413.4\) & \(80.3625\) &
\(91.1879\) & \(1\) \\*
second & KS & \(0.111572\) & \(2206.8\) & \(80.3625\) & \(91.1879\) &
\(1\) \\*
second & CJT & \(0.111572\) & \(2206.8\) & \(80.3625\) & \(91.1879\) &
\(1\) \\*
first & HS (1, 1) & \(0.0816903\) & \(3014.1\) & \(80.3625\) &
\(81.6624\) & \(1.24689\) \\*
first & HS (0, 1) & \(0.0799859\) & \(3078.3\) & \(80.3625\) &
\(82.0733\) & \(1.23444\) \\*
first & HS (0, 0) & \(0.0799859\) & \(3078.3\) & \(80.3625\) &
\(82.244\) & \(1.22932\) \\*
first & KS & \(0.0868678\) & \(2834.4\) & \(80.3625\) & \(83.1831\) &
\(1.20172\) \\*
first & CJT & \(0.0868678\) & \(2834.4\) & \(80.3625\) & \(83.1831\) &
\(1.20172\) \\
\end{longtable}

The masses of the fermionic excitations follow from stage 3. In the
second version every constituent is paired, and the lowest
quasi-particle energy, the gap \(E^{X,{\rm gap}}_{\rm SH}\), belongs to
a doublet of \(SU(2)_{\rm V}\) with the electric charges \(\tfrac34\)
and \(-\tfrac14\). In GeV it is
\(E^{X,{\rm gap}}_{\rm SH}\Lambda_{\rm LL}\):

Table 4.2. The gap of the ground states (stage 3) in units of the cutoff
and in GeV.

\sptkeptpenalty=10000
\begin{longtable}[]{@{}
  >{\raggedright\arraybackslash}p{(\columnwidth - 6\tabcolsep) * \real{0.1818}}
  >{\raggedright\arraybackslash}p{(\columnwidth - 6\tabcolsep) * \real{0.2727}}
  >{\raggedright\arraybackslash}p{(\columnwidth - 6\tabcolsep) * \real{0.2727}}
  >{\raggedright\arraybackslash}p{(\columnwidth - 6\tabcolsep) * \real{0.2727}}@{}}
\toprule\noalign{}
\begin{minipage}[b]{\linewidth}\raggedright
version
\end{minipage} & \begin{minipage}[b]{\linewidth}\raggedright
method
\end{minipage} & \begin{minipage}[b]{\linewidth}\raggedright
\(E^{X,{\rm gap}}_{\rm SH}\) (units of \(\Lambda_{\rm LL}\))
\end{minipage} & \begin{minipage}[b]{\linewidth}\raggedright
\(E^{X,{\rm gap}}_{\rm SH}\) (GeV)
\end{minipage} \\
\midrule\noalign{}
\endhead
\bottomrule\noalign{}
\endlastfoot
second & HS (1, 1) & \(0.426792\) & \(1004.4\) \\*
second & HS (0, 1) & \(0.452967\) & \(1093.2\) \\*
second & HS (0, 0) & \(0.452967\) & \(1093.2\) \\*
second & KS & \(0.384201\) & \(847.86\) \\*
second & CJT & \(0.384201\) & \(847.86\) \\*
first & HS (1, 1) & \(0.170392\) & \(513.57\) \\*
first & HS (0, 1) & \(0.2\) & \(615.66\) \\*
first & HS (0, 0) & \(0.2\) & \(615.66\) \\*
first & KS & \(0.111459\) & \(315.92\) \\*
first & CJT & \(0.111459\) & \(315.92\) \\
\end{longtable}

\hypertarget{the-amplitude-mode}{%
\section{5. The amplitude mode}\label{the-amplitude-mode}}

\hypertarget{the-effective-lagrangian}{%
\subsection{5.1 The effective
Lagrangian}\label{the-effective-lagrangian}}

At long wavelengths the pair fields obey the effective Lagrangian
\(Z_{\rm SH}\big(\lvert\partial\Phi^{\rm d}\rvert^2+\lvert\partial\Phi'^{\rm u}\rvert^2\big)-V^{\rm eff}_{\rm SH}\)
(second version), where \(\Phi^{\rm d}\) and \(\Phi'^{\rm u}\) are the
neutral components \(\langle\Pi^{\rm d}_{\rm SH}\rangle\) and
\(\langle\Pi'^{\rm u}_{\rm SH}\rangle\). The kinetic coefficient follows
from the phase stiffness. A chemical potential \(\mu\) of
\(T^{3,{\rm SM}}\), under which \(\Phi^{\rm d}\) has the charge
\(-\tfrac12\) and \(\Phi'^{\rm u}\) the charge \(+\tfrac12\), lowers the
energy by \(Z_{\rm SH}\mu^2(\tfrac14+\tfrac14)(\varphi^0_{\rm SH})^2\),
so \(\chi^{33}_{\rm SH}=Z_{\rm SH}(\varphi^0_{\rm SH})^2\). When both
lengths oscillate together,
\(\varphi_{\rm SH}=\varphi^0_{\rm SH}+\eta\), the kinetic term is
\(2Z_{\rm SH}\dot\eta^2\) and the potential
\(\tfrac12V^{\rm eff\,\prime\prime}_{\rm SH}\eta^2\). The frequency of
the oscillation is
\[\big(M^{\rm EFT}_{\rm SH}\big)^2=\frac{V^{\rm eff\,\prime\prime}_{\rm SH}(\varphi^0_{\rm SH})\,(\varphi^0_{\rm SH})^2}{4\chi^{33}_{\rm SH}}.\tag{5.1}\]
The same formula holds in the first version, with one pair field: there
\(\chi^{33}_{\rm SH}=\tfrac12Z_{\rm SH}(\varphi^0_{\rm SH})^2\) and the
kinetic term is \(Z_{\rm SH}\dot\eta^2\) (check WH-2, symbolic). Formula
(5.1) is the leading order of an expansion in the frequency. It is
reliable when the mode lies well below the energy at which a pair can
break into two quasi-particles, which is \(2E^{X,{\rm gap}}_{\rm SH}\)
in the second version, where every field is paired. Here it does not, as
Table 5.1 shows.

\hypertarget{the-linear-response}{%
\subsection{5.2 The linear response}\label{the-linear-response}}

The full answer at zero momentum is the pole of the linear response of
the time-dependent mean field. A uniform fluctuation with the time
factor \(e^{-i\omega t}\) changes the self-energies by complex
amplitudes:

\begin{itemize}
\tightlist
\item
  \(\delta\mathsf N\), the change of the normal self-energy;
\item
  \(\delta\mathsf P\) and \(\delta\bar{\mathsf P}\), the changes of the
  pair self-energy and of its adjoint (at a real frequency these are
  independent amplitudes).
\end{itemize}

These change the local matrices by \(\boldsymbol\chi_{\rm SH}(\omega)\),
the linear response of the state at the frequency \(\omega\). In the
eigenbasis \((E_k,\lvert k\rangle)\) of (1.3) at each momentum it is
\[\delta R_{kl}=\delta H_{kl}\,\frac{f_l-f_k}{\omega-E_k+E_l},\qquad f_k=1\text{ for }E_k<0,\ 0\text{ otherwise},\tag{5.2}\]
summed over the momenta, with \(\delta H\) the change of (1.3) and
\(\delta R\) that of the projector onto its negative eigenvalues. The
contact terms map the changes of the local matrices back to changes of
the self-energies by the kernel \(\mathsf F^X_{\rm SH}\) of method
\(X\). A collective mode is a frequency at which
\(1-\boldsymbol\chi_{\rm SH}(\omega)\mathsf F^X_{\rm SH}\) is singular.
At \(\omega=0\) this is the static response of stage 3. A broken
symmetry gives a zero there: the Goldstone mode.

The program computes
\(1-\boldsymbol\chi_{\rm SH}(\omega)\mathsf F^X_{\rm SH}\) in the sector
of relative number 0: the normal blocks of \(\Psi_{\rm LL}\) and of the
pair of singlets and the pairs of \(\Psi_{\rm LL}\) with the singlets,
closed under the symmetries of the state. It does so in HS (0, 1) and HS
(0, 0), where the amplitude mode lies below the threshold of two paired
quasi-particles, \(2M^{\rm eff}_{\rm SH}\) with
\(M^{\rm eff}_{\rm SH}=(m_{\rm LL}^2+(d^0_{\rm SH})^2/4)^{1/2}\), and is
a sharp state. It searches the frequencies below the threshold for zeros
of the smallest singular value (600 equal steps, and 25 more that
approach the threshold geometrically, refined by golden sections). The
eigenvalues of
\((1-\boldsymbol\chi_{\rm SH}\mathsf F^X_{\rm SH})^\dagger(1-\boldsymbol\chi_{\rm SH}\mathsf F^X_{\rm SH})\)
give the singular values only to about the square root of the machine
precision, which would fix a zero only to about \(10^{-8}\). The program
therefore takes the singular values of
\(1-\boldsymbol\chi_{\rm SH}\mathsf F^X_{\rm SH}\) restricted to the
eigenvectors of that product whose eigenvalues lie below \(10^{-6}\),
which keep the precision of the machine; at the modes found they are
below \(2\times10^{-15}\). At each zero it finds the null space, maps it
to the field amplitudes by \(\mathsf F^X_{\rm SH}\), and projects onto
it the amplitude of the vacuum direction, the fluctuation
\((\delta\mathsf P,\delta\bar{\mathsf P})=(\mathsf P^{\rm vac}_{\rm SH},(\mathsf P^{\rm vac}_{\rm SH})^\dagger)\),
normalized. The amplitude mode is the zero onto whose null space that
projection is (nearly) the whole vector; the dimension of the null space
is its multiplicity. Its frequency is written
\(\omega^{\rm H}_{\rm SH}\).

The results, in both versions and both methods (checks H-7a to H-7e):

\begin{itemize}
\tightlist
\item
  at \(\omega=0\) exactly 4 (second version) and 3 (first version)
  singular values vanish: the Goldstone modes of the broken generators;
\item
  below the threshold there is one zero whose null space contains the
  amplitude of the vacuum direction; its multiplicity is 4 in the second
  version and 1 in the first;
\item
  the frequency of that zero equals \(d^0_{\rm SH}\) to
  \(2\times10^{-13}\) (relative; Section 5.3 shows why);
\item
  HS (0, 1) and HS (0, 0) give the same frequency, because the vector
  term of HS (0, 1) acts on the number density and the currents, which
  the amplitude does not change;
\item
  the lowest transition energy whose terms in
  \(\boldsymbol\chi_{\rm SH}(\omega)\) act on the null space is the
  threshold of two paired quasi-particles, \(2M^{\rm eff}_{\rm SH}\) (up
  to the smallest momentum of the grid). In the first version the sector
  also holds the transitions of the unpaired fields \(\Psi^{\rm u}\) and
  \(\chi'_{\rm SH}\), from \(0.400029\) on. Their terms vanish on the
  null space (the part of each that acts there is below
  \(2\times10^{-31}\) of its size, the level of the rounding), so they
  do not act on the mode, which stays sharp although it lies above them.
\end{itemize}

The Python program computes the amplitude mode of HS (0, 0) a second
way, from the response of the pair amplitudes alone. In that method the
kernel reads only those amplitudes. The phase eigenvalue of its response
is 1 at \(\omega=0\), a check of the normalization, and the amplitude
and the phase do not couple (check YH-4).

\hypertarget{the-closed-form-of-the-pairing-only-methods-1}{%
\subsection{5.3 The closed form of the pairing-only
methods}\label{the-closed-form-of-the-pairing-only-methods-1}}

In HS (0, 0) the response of the pair amplitudes takes the form that the
theory of Bardeen, Cooper and Schrieffer gives for a superconductor.
With the energy \(E=(q^2+(M^{\rm eff}_{\rm SH})^2)^{1/2}\) of a
quasi-particle of momentum \(q\) (Section 2.2), the eigenvalues along
the amplitude and along the phase of the vacuum direction are
\[\lambda^{\rm amp}_{\rm SH}(\omega)=\frac{G_{\rm SH}}{4\pi^2}\int_0^{\Lambda_{\rm LL}}\frac{q^2\,(q^2+m_{\rm LL}^2)\,dq}{E\,(E^2-\omega^2/4)},\qquad
\lambda^{\rm ph}_{\rm SH}(\omega)=\frac{G_{\rm SH}}{4\pi^2}\int_0^{\Lambda_{\rm LL}}\frac{q^2\,E\,dq}{E^2-\omega^2/4}.\tag{5.3}\]
The programs evaluate these integrals as sums over the nodes of their
grid, with the weights of the grid. The Python program finds that its
response equals the sums (5.3) to \(10^{-15}\) at five frequencies, in
both versions (check YH-7). I found these forms by that comparison; I
have not derived them from (5.2) by hand.

Two consequences follow.

\begin{itemize}
\tightlist
\item
  At \(\omega=0\) the phase eigenvalue is
  \(\lambda^{\rm ph}_{\rm SH}(0)=G_{\rm SH}I^{(1)}_{\rm SF}(M^{\rm eff}_{\rm SH})/(4\pi^2)=1\).
  This is the gap equation of Section 2.2, and it is the Goldstone
  theorem: the phase of the condensate costs no energy.
\item
  Since \(E^2=q^2+m_{\rm LL}^2+(d^0_{\rm SH})^2/4\), the integrands
  obey, at every momentum,
  \[\frac{q^2+m_{\rm LL}^2}{E\,(E^2-\omega^2/4)}-\frac1E=\frac{\omega^2-(d^0_{\rm SH})^2}{4E\,(E^2-\omega^2/4)}.\]
  With the gap equation this gives
  \[\lambda^{\rm amp}_{\rm SH}(\omega)-1=\frac{\omega^2-(d^0_{\rm SH})^2}{4}\,\frac{G_{\rm SH}}{4\pi^2}\int_0^{\Lambda_{\rm LL}}\frac{q^2\,dq}{E\,(E^2-\omega^2/4)}.\tag{5.4}\]
  Below the threshold, \(\omega<2M^{\rm eff}_{\rm SH}\), every term of
  the integral is positive. So \(\lambda^{\rm amp}_{\rm SH}(\omega)=1\)
  only at \[\omega^{\rm H}_{\rm SH}=d^0_{\rm SH},\tag{5.5}\] exactly,
  for every coupling, mass and cutoff, and on every grid with positive
  weights. The Wolfram script derives the identity and the positivity
  symbolically and finds the root of \(\lambda^{\rm amp}_{\rm SH}=1\) at
  \(d^0_{\rm SH}\) to 22 digits (check WH-5).
\end{itemize}

The Rust program, which solves the linear response of the whole sector
and does not use (5.3), finds (5.5) to \(2\times10^{-13}\) (check H-7d)
and in all 32 points of the scan of Section 6 to \(2\times10^{-12}\)
(check H-8d).

Result (5.5) is the analogue of a known result of the theory of Bardeen,
Cooper and Schrieffer: the amplitude mode of a superconductor lies at
twice the gap parameter. Here \(d^0_{\rm SH}/2\) plays the role of the
gap parameter. In a superconductor the lowest energy of two
quasi-particles is also twice the gap parameter, so the mode lies at the
threshold. Here the mass \(m_{\rm LL}\) of the fields raises the lowest
energy of a quasi-particle from \(d^0_{\rm SH}/2\) to
\(M^{\rm eff}_{\rm SH}\). The mode is therefore bound, by
\(2M^{\rm eff}_{\rm SH}-d^0_{\rm SH}=(4m_{\rm LL}^2+(d^0_{\rm SH})^2)^{1/2}-d^0_{\rm SH}=0.0930893\)
for the couplings of stage 3.

\textbf{The partners of the amplitude mode in the second version.} The
four complex pair fields \(\langle\Pi^{\rm u}_{\rm SH}\rangle\),
\(\langle\Pi^{\rm d}_{\rm SH}\rangle\),
\(\langle\Pi'^{\rm u}_{\rm SH}\rangle\),
\(\langle\Pi'^{\rm d}_{\rm SH}\rangle\) have eight real components. On
all of them the Python program finds four eigenvalues equal to 1 at
\(\omega=0\) and four more at \(\omega=d^0_{\rm SH}\) (check YH-8). The
generators \(T^{a,{\rm SM}}+T^{a,{\rm R}}_{\rm SH}\) of the custodial
group \(SU(2)_{\rm V}\) commute with the response, and on each set of
four modes the Casimir
\(\sum_a(T^{a,{\rm SM}}+T^{a,{\rm R}}_{\rm SH})^2\) has the eigenvalues
0, 2, 2, 2: a singlet and a triplet. The modes are:

\begin{itemize}
\tightlist
\item
  at \(\omega=0\), the triplet of Goldstone modes that become the
  longitudinal parts of \(W^\pm_\mu\) and \(Z_\mu\), and the singlet
  Goldstone mode of the particle number \(F^{\rm tot}_{\rm SH}\)
  (decision D32 of the author);
\item
  at \(\omega=d^0_{\rm SH}\), the singlet amplitude mode of the vacuum
  direction, the candidate for the Higgs boson, and a triplet of three
  more massive scalars. The electric charge \(Q^{\rm SM}\) equals
  \(T^{3,{\rm SM}}+T^{3,{\rm R}}_{\rm SH}\) on the pair fields, so the
  triplet holds a neutral scalar and a charged scalar with its
  antiparticle.
\end{itemize}

The second version therefore has the scalar content of two Higgs
doublets, not of one. In the first version, with one doublet, the same
count gives three Goldstone modes and one amplitude mode, as in the
Standard Model (check YH-8). The multiplicity of the amplitude mode in
the Rust program, 4 in the second version and 1 in the first, is the
same count. The degeneracy of the triplet with the singlet is a property
of the pairing-only methods. The dynamics in time of stage 6, which
follows the time-dependent mean field of every method, finds that the
scalar and the vector terms split it: in HS (1, 1) of the second version
the triplet oscillates at 0.828525, 2.2 per cent above the singlet at
0.810871.

\hypertarget{the-results-1}{%
\subsection{5.4 The results}\label{the-results-1}}

Table 5.1. The amplitude mode from the effective Lagrangian (all
methods), in units of the cutoff, relative to \(v_{\rm SH}\) and in GeV,
and twice the gap \(E^{X,{\rm gap}}_{\rm SH}\) of stage 3, the lowest
energy of two quasi-particles. In the second version every field is
paired and \(2E^{X,{\rm gap}}_{\rm SH}\) is the energy that breaks a
pair; in the first version the gap is that of the unpaired fields.

\sptkeptpenalty=10000
\begin{longtable}[]{@{}
  >{\raggedright\arraybackslash}p{(\columnwidth - 10\tabcolsep) * \real{0.1321}}
  >{\raggedright\arraybackslash}p{(\columnwidth - 10\tabcolsep) * \real{0.1887}}
  >{\raggedright\arraybackslash}p{(\columnwidth - 10\tabcolsep) * \real{0.1698}}
  >{\raggedright\arraybackslash}p{(\columnwidth - 10\tabcolsep) * \real{0.1698}}
  >{\raggedright\arraybackslash}p{(\columnwidth - 10\tabcolsep) * \real{0.1698}}
  >{\raggedright\arraybackslash}p{(\columnwidth - 10\tabcolsep) * \real{0.1698}}@{}}
\toprule\noalign{}
\begin{minipage}[b]{\linewidth}\raggedright
version
\end{minipage} & \begin{minipage}[b]{\linewidth}\raggedright
method
\end{minipage} & \begin{minipage}[b]{\linewidth}\raggedright
\(M^{\rm EFT}_{\rm SH}\)
\end{minipage} & \begin{minipage}[b]{\linewidth}\raggedright
\(M^{\rm EFT}_{\rm SH}/v_{\rm SH}\)
\end{minipage} & \begin{minipage}[b]{\linewidth}\raggedright
\(M^{\rm EFT}_{\rm SH}\) (GeV)
\end{minipage} & \begin{minipage}[b]{\linewidth}\raggedright
\(2E^{X,{\rm gap}}_{\rm SH}\)
\end{minipage} \\
\midrule\noalign{}
\endhead
\bottomrule\noalign{}
\endlastfoot
second & HS (1, 1) & \(0.96457\) & \(9.219\) & \(2270\) & \(0.85358\) \\*
second & HS (0, 1) & \(0.94638\) & \(9.276\) & \(2284\) & \(0.90593\) \\*
second & HS (0, 0) & \(0.94638\) & \(9.276\) & \(2284\) & \(0.90593\) \\*
second & KS & \(1.0676\) & \(9.569\) & \(2356\) & \(0.7684\) \\*
second & CJT & \(1.0676\) & \(9.569\) & \(2356\) & \(0.7684\) \\*
first & HS (1, 1) & \(0.95778\) & \(11.72\) & \(2887\) & \(0.34078\) \\*
first & HS (0, 1) & \(0.94835\) & \(11.86\) & \(2919\) & \(0.4\) \\*
first & HS (0, 0) & \(0.94638\) & \(11.83\) & \(2913\) & \(0.4\) \\*
first & KS & \(0.97123\) & \(11.18\) & \(2753\) & \(0.22292\) \\*
first & CJT & \(0.97123\) & \(11.18\) & \(2753\) & \(0.22292\) \\
\end{longtable}

Table 5.2. The amplitude mode from the linear response (HS (0, 1), HS
(0, 0)): its frequency, its relative difference from \(d^0_{\rm SH}\),
its ratio to \(v_{\rm SH}\), its value in GeV, the lowest transition
energy that acts on it, and its multiplicity.

\sptkeptpenalty=10000
\begin{longtable}[]{@{}
  >{\raggedright\arraybackslash}p{(\columnwidth - 14\tabcolsep) * \real{0.0889}}
  >{\raggedright\arraybackslash}p{(\columnwidth - 14\tabcolsep) * \real{0.1222}}
  >{\raggedright\arraybackslash}p{(\columnwidth - 14\tabcolsep) * \real{0.1667}}
  >{\raggedright\arraybackslash}p{(\columnwidth - 14\tabcolsep) * \real{0.1444}}
  >{\raggedright\arraybackslash}p{(\columnwidth - 14\tabcolsep) * \real{0.1000}}
  >{\raggedright\arraybackslash}p{(\columnwidth - 14\tabcolsep) * \real{0.0889}}
  >{\raggedright\arraybackslash}p{(\columnwidth - 14\tabcolsep) * \real{0.1667}}
  >{\raggedright\arraybackslash}p{(\columnwidth - 14\tabcolsep) * \real{0.1222}}@{}}
\toprule\noalign{}
\begin{minipage}[b]{\linewidth}\raggedright
version
\end{minipage} & \begin{minipage}[b]{\linewidth}\raggedright
method
\end{minipage} & \begin{minipage}[b]{\linewidth}\raggedright
\(\omega^{\rm H}_{\rm SH}\)
\end{minipage} & \begin{minipage}[b]{\linewidth}\raggedright
\(\omega^{\rm H}_{\rm SH}/d^0_{\rm SH}-1\)
\end{minipage} & \begin{minipage}[b]{\linewidth}\raggedright
\(\omega^{\rm H}_{\rm SH}/v_{\rm SH}\)
\end{minipage} & \begin{minipage}[b]{\linewidth}\raggedright
\(\omega^{\rm H}_{\rm SH}\) (GeV)
\end{minipage} & \begin{minipage}[b]{\linewidth}\raggedright
lowest transition that acts on it
\end{minipage} & \begin{minipage}[b]{\linewidth}\raggedright
multiplicity
\end{minipage} \\
\midrule\noalign{}
\endhead
\bottomrule\noalign{}
\endlastfoot
second & HS (0, 1) & \(0.8128455463\) & \(-2\times10^{-13}\) & \(7.967\)
& \(1962\) & \(0.905948\) & 4 \\*
second & HS (0, 0) & \(0.8128455463\) & \(-2\times10^{-13}\) & \(7.967\)
& \(1962\) & \(0.905948\) & 4 \\*
first & HS (0, 1) & \(0.8128455463\) & \(-2\times10^{-13}\) & \(10.16\)
& \(2502\) & \(0.905948\) & 1 \\*
first & HS (0, 0) & \(0.8128455463\) & \(-2\times10^{-13}\) & \(10.16\)
& \(2502\) & \(0.905948\) & 1 \\
\end{longtable}

\textbf{What the tables show.}

\begin{itemize}
\tightlist
\item
  \textbf{The amplitude mode lies at
  \(\omega^{\rm H}_{\rm SH}=d^0_{\rm SH}\).} In HS (0, 1) and HS (0, 0)
  the linear response puts it at \(\omega^{\rm H}_{\rm SH}=0.812846\) in
  both versions, which is \(d^0_{\rm SH}\) (Section 5.3), below the
  threshold \(2M^{\rm eff}_{\rm SH}=0.905935\).
\item
  \textbf{The effective Lagrangian fails here.} In the second version it
  puts the mode at \(0.9464\) to \(1.068\), above the energy
  \(2E^{X,{\rm gap}}_{\rm SH}\) that breaks a pair in every method,
  outside its range of validity. In HS (0, 1) and HS (0, 0) it
  overestimates the frequency of the linear response by the factor
  \(1.164\).
\item
  \textbf{The second version lowers the ratio, but not enough.} At the
  same frequency of the amplitude mode, the second doublet increases
  \(v_{\rm SH}\) by the factor \(1.276\) (HS (0, 0)). So the ratio
  \(M_{\rm H}/v_{\rm SH}\) falls from \(10.16\) (first version) to
  \(7.967\) (second version).
\item
  \textbf{The comparison with the Standard Model.} There
  \(m_{\rm H}/v=0.5082\) (\(m_{\rm H}=125.13\) GeV). The amplitude mode
  of the superfluid is \(16\) to \(19\) times heavier relative to the
  scale, about \(2\times10^{3}\) to \(2.4\times10^{3}\) GeV.
\end{itemize}

\hypertarget{the-dependence-on-the-coupling-and-the-mass}{%
\section{6. The dependence on the coupling and the
mass}\label{the-dependence-on-the-coupling-and-the-mass}}

The coupling \(G_{\rm SH}\) was a choice of stage 3,
\(G_{\rm SH}/G^{\rm crit}_{\rm SH}=1.1838\) with
\(G^{\rm crit}_{\rm SH}=4\pi^2/I^{(1)}_{\rm SF}(m_{\rm LL})\), and so
was the mass. In HS (0, 0) the ground state has the closed form of
Section 2.2 for any coupling and mass. The program therefore repeats the
computation of \(v_{\rm SH}\), \(\rho_{\rm SH}\) and the amplitude mode
for \(G_{\rm SH}/G^{\rm crit}_{\rm SH}\) from 1.02 to 3 and the masses
0.1 and 0.2 of the three fields, in both versions. The amplitude mode is
computed in the sector of the neutral pairs of the vacuum alone, which
is exact in this method (check H-8c compares with the full sector).

Table 6.1. The scan (HS (0, 0)), for the common mass
\(m_{\rm LL}=m^\chi_{\rm SH}=m^{\chi'}_{\rm SH}\) of the three fields:
the ratio \(M_{\rm H}/v_{\rm SH}=\omega^{\rm H}_{\rm SH}/v_{\rm SH}\)
and \(\rho_{\rm SH}\).

\begin{longtable}[]{@{}
  >{\raggedright\arraybackslash}p{(\columnwidth - 16\tabcolsep) * \real{0.0952}}
  >{\raggedright\arraybackslash}p{(\columnwidth - 16\tabcolsep) * \real{0.0714}}
  >{\raggedright\arraybackslash}p{(\columnwidth - 16\tabcolsep) * \real{0.1190}}
  >{\raggedright\arraybackslash}p{(\columnwidth - 16\tabcolsep) * \real{0.1190}}
  >{\raggedright\arraybackslash}p{(\columnwidth - 16\tabcolsep) * \real{0.1190}}
  >{\raggedright\arraybackslash}p{(\columnwidth - 16\tabcolsep) * \real{0.1190}}
  >{\raggedright\arraybackslash}p{(\columnwidth - 16\tabcolsep) * \real{0.1190}}
  >{\raggedright\arraybackslash}p{(\columnwidth - 16\tabcolsep) * \real{0.1190}}
  >{\raggedright\arraybackslash}p{(\columnwidth - 16\tabcolsep) * \real{0.1190}}@{}}
\toprule\noalign{}
\begin{minipage}[b]{\linewidth}\raggedright
version
\end{minipage} & \begin{minipage}[b]{\linewidth}\raggedright
\(m_{\rm LL}\)
\end{minipage} & \begin{minipage}[b]{\linewidth}\raggedright
\(G_{\rm SH}/G^{\rm crit}_{\rm SH}\)
\end{minipage} & \begin{minipage}[b]{\linewidth}\raggedright
\(d^0_{\rm SH}\)
\end{minipage} & \begin{minipage}[b]{\linewidth}\raggedright
\(v_{\rm SH}\)
\end{minipage} & \begin{minipage}[b]{\linewidth}\raggedright
\(\rho_{\rm SH}\)
\end{minipage} & \begin{minipage}[b]{\linewidth}\raggedright
\(\omega^{\rm H}_{\rm SH}\)
\end{minipage} & \begin{minipage}[b]{\linewidth}\raggedright
\(2M^{\rm eff}_{\rm SH}\)
\end{minipage} & \begin{minipage}[b]{\linewidth}\raggedright
\(\omega^{\rm H}_{\rm SH}/v_{\rm SH}\)
\end{minipage} \\
\midrule\noalign{}
\endhead
\bottomrule\noalign{}
\endlastfoot
second & \(0.1\) & \(1.02\) & \(0.20583\) & \(0.042078\) & \(1\) &
\(0.205833\) & \(0.286997\) & \(4.892\) \\
second & \(0.1\) & \(1.05\) & \(0.34088\) & \(0.062877\) & \(1\) &
\(0.340883\) & \(0.395223\) & \(5.421\) \\
second & \(0.1\) & \(1.1\) & \(0.50977\) & \(0.082828\) & \(1\) &
\(0.509767\) & \(0.547597\) & \(6.155\) \\
second & \(0.1\) & \(1.1838\) & \(0.73828\) & \(0.1021\) & \(1\) &
\(0.738278\) & \(0.764889\) & \(7.231\) \\
second & \(0.1\) & \(1.3\) & \(1.0064\) & \(0.11671\) & \(1\) &
\(1.00635\) & \(1.02604\) & \(8.623\) \\
second & \(0.1\) & \(1.5\) & \(1.4071\) & \(0.12807\) & \(1\) &
\(1.40709\) & \(1.42123\) & \(10.99\) \\
second & \(0.1\) & \(2\) & \(2.2705\) & \(0.13152\) & \(1\) & \(2.2705\)
& \(2.2793\) & \(17.26\) \\
second & \(0.1\) & \(3\) & \(3.8001\) & \(0.11903\) & \(1\) &
\(3.80008\) & \(3.80534\) & \(31.92\) \\
second & \(0.2\) & \(1.02\) & \(0.24066\) & \(0.041741\) & \(1\) &
\(0.240656\) & \(0.466814\) & \(5.765\) \\
second & \(0.2\) & \(1.05\) & \(0.39029\) & \(0.062858\) & \(1\) &
\(0.390293\) & \(0.558864\) & \(6.209\) \\
second & \(0.2\) & \(1.1\) & \(0.57181\) & \(0.082962\) & \(1\) &
\(0.571812\) & \(0.697832\) & \(6.892\) \\
second & \(0.2\) & \(1.1838\) & \(0.81285\) & \(0.10202\) & \(1\) &
\(0.812846\) & \(0.905935\) & \(7.967\) \\
second & \(0.2\) & \(1.3\) & \(1.0931\) & \(0.11615\) & \(1\) &
\(1.09305\) & \(1.16394\) & \(9.411\) \\
second & \(0.2\) & \(1.5\) & \(1.5106\) & \(0.12679\) & \(1\) &
\(1.51061\) & \(1.56267\) & \(11.91\) \\
second & \(0.2\) & \(2\) & \(2.411\) & \(0.12934\) & \(1\) & \(2.41104\)
& \(2.444\) & \(18.64\) \\
second & \(0.2\) & \(3\) & \(4.0115\) & \(0.11652\) & \(1\) &
\(4.01146\) & \(4.03136\) & \(34.43\) \\
first & \(0.1\) & \(1.02\) & \(0.20583\) & \(0.030618\) & \(1.05896\) &
\(0.205833\) & \(0.286997\) & \(6.723\) \\
first & \(0.1\) & \(1.05\) & \(0.34088\) & \(0.046903\) & \(1.11287\) &
\(0.340883\) & \(0.395223\) & \(7.268\) \\
first & \(0.1\) & \(1.1\) & \(0.50977\) & \(0.06347\) & \(1.1744\) &
\(0.509767\) & \(0.547597\) & \(8.032\) \\
first & \(0.1\) & \(1.1838\) & \(0.73828\) & \(0.080797\) & \(1.25253\)
& \(0.738278\) & \(0.764889\) & \(9.137\) \\
first & \(0.1\) & \(1.3\) & \(1.0064\) & \(0.095641\) & \(1.34313\) &
\(1.00635\) & \(1.02604\) & \(10.52\) \\
first & \(0.1\) & \(1.5\) & \(1.4071\) & \(0.11017\) & \(1.48008\) &
\(1.40709\) & \(1.42123\) & \(12.77\) \\
first & \(0.1\) & \(2\) & \(2.2705\) & \(0.12369\) & \(1.76891\) &
\(2.2705\) & \(2.2793\) & \(18.36\) \\
first & \(0.1\) & \(3\) & \(3.8001\) & \(0.12476\) & \(2.19689\) &
\(3.80008\) & \(3.80534\) & \(30.46\) \\
first & \(0.2\) & \(1.02\) & \(0.24066\) & \(0.030059\) & \(1.03721\) &
\(0.240656\) & \(0.466814\) & \(8.006\) \\
first & \(0.2\) & \(1.05\) & \(0.39029\) & \(0.046256\) & \(1.08303\) &
\(0.390293\) & \(0.558864\) & \(8.438\) \\
first & \(0.2\) & \(1.1\) & \(0.57181\) & \(0.062771\) & \(1.14496\) &
\(0.571812\) & \(0.697832\) & \(9.109\) \\
first & \(0.2\) & \(1.1838\) & \(0.81285\) & \(0.079986\) & \(1.22932\)
& \(0.812846\) & \(0.905935\) & \(10.16\) \\
first & \(0.2\) & \(1.3\) & \(1.0931\) & \(0.09464\) & \(1.32786\) &
\(1.09305\) & \(1.16394\) & \(11.55\) \\
first & \(0.2\) & \(1.5\) & \(1.5106\) & \(0.10886\) & \(1.47434\) &
\(1.51061\) & \(1.56267\) & \(13.88\) \\
first & \(0.2\) & \(2\) & \(2.411\) & \(0.12187\) & \(1.77559\) &
\(2.41104\) & \(2.444\) & \(19.78\) \\
first & \(0.2\) & \(3\) & \(4.0115\) & \(0.12254\) & \(2.21203\) &
\(4.01146\) & \(4.03136\) & \(32.74\) \\
\end{longtable}

In every point the amplitude mode lies at
\(\omega^{\rm H}_{\rm SH}=d^0_{\rm SH}\) (check H-8d), as (5.5)
requires, so the ratio is \(d^0_{\rm SH}/v_{\rm SH}\). It stays between
\(4.892\) and \(34.43\) in the second version and between \(6.723\) and
\(32.74\) in the first. Near the critical coupling both \(v_{\rm SH}\)
and the amplitude mode shrink with the condensate, and their ratio does
not approach the value \(0.508\) of the Standard Model.

\hypertarget{what-the-results-mean}{%
\section{7. What the results mean}\label{what-the-results-mean}}

\begin{itemize}
\tightlist
\item
  \textbf{What works.} The superfluid of the second version has the
  quantum numbers of the Higgs field, breaks \(SU(2)_L\times U(1)_Y\) to
  the \(U(1)\) of the electric charge, leaves the photon massless and
  gives \(\rho_{\rm SH}=1\). With the scale set to \(v=246.22\) GeV, the
  \(W\) and the \(Z\) have the masses of the Standard Model at tree
  level.
\item
  \textbf{What fails.} The amplitude mode, the candidate for the Higgs
  boson, is \(16\) to \(19\) times too heavy relative to the scale. In a
  superfluid of fermion pairs bound by a contact interaction the
  amplitude mode lies at \(d^0_{\rm SH}\), just below the energy
  \(2M^{\rm eff}_{\rm SH}\) that breaks a pair (Section 5.3), while
  \(v_{\rm SH}^2\), the phase stiffness, grows only like
  \((M^{\rm eff}_{\rm SH})^2\) times the logarithm of
  \(\Lambda_{\rm LL}^2/(M^{\rm eff}_{\rm SH})^2\) (in the leading
  approximation, divided by \(4\pi^2\) and multiplied by the number of
  paired species). Here the cutoff is only about twice
  \(M^{\rm eff}_{\rm SH}\), so \(v_{\rm SH}\) is small compared with
  \(M^{\rm eff}_{\rm SH}\) and the ratio is large. The same relation
  underlies the models in which the Higgs doublet is a bound state of
  the top quark and its antiparticle, where a light Higgs boson needs a
  very high cutoff (Bardeen, Hill and Lindner 1990).
\item
  \textbf{What the second version adds.} Its two pair fields carry the
  scalars of two Higgs doublets: besides the candidate for the Higgs
  boson, a custodial triplet of massive scalars (one neutral, one
  charged with its antiparticle), degenerate with it in the pairing-only
  methods, and the extra Goldstone boson of the particle number (Section
  5.3).
\item
  \textbf{What remains open.} The mean field leaves out the fluctuations
  of the pair fields, which shift both the stiffness and the amplitude
  mode; the extra Goldstone boson (decision D32 of the author); and the
  fermions with fractional charges at about \(848\) to
  \(1.09\times10^{3}\) GeV.
\end{itemize}

\hypertarget{the-programs-and-the-checks}{%
\section{8. The programs and the
checks}\label{the-programs-and-the-checks}}

The results were computed and checked by three programs written for this
stage:

\begin{itemize}
\tightlist
\item
  \textbf{Rust.} The program \texttt{sh\_higgs} (crate
  \texttt{rust/shhiggs}), built on the engine \texttt{uniform\_n} and
  the module \texttt{rpa\_n} of the library \texttt{sfcore} (the linear
  response at a frequency of the states of several Dirac fields). It
  includes the model of stage 3 unchanged. It ran its 27 checks in 2174
  s on two parallel threads; the longest of the ten computations (five
  methods, two versions) took 411 s. Its verdict: ALL PASS (27 of 27).
\item
  \textbf{Python.} The program \texttt{python/sh\_higgs.py} uses the
  engine of the Python program of stage 3. It computes the
  susceptibilities from the energy, the curvature of HS (0, 0) and the
  amplitude mode by the reduced linear response, compares that response
  with the closed forms (5.3), counts the modes on all the pair fields
  with their custodial content, and compares its numbers with the data
  of the Rust program: ALL PASS (9 of 9).
\item
  \textbf{Wolfram.} The script \texttt{wolfram/source/sh\_higgs.wls},
  run on the author's computer, derives the custodial algebra, the
  formula of the amplitude mode of the effective Lagrangian and the
  identity behind (5.5) symbolically, evaluates the closed forms of HS
  (0, 0) to 30 digits and checks the conversion to GeV: ALL PASS (5 of
  5).
\end{itemize}

The reports are \texttt{data/hck\_sh\_higgs\_checks.txt},
\texttt{data/hck\_sh\_higgs\_checks\_python.txt} and
\texttt{data/hck\_sh\_higgs\_checks\_wolfram.txt}. The data are:

\begin{itemize}
\tightlist
\item
  \texttt{data/sh\_hg\_potential.csv}: the curves of the effective
  potential;
\item
  \texttt{data/sh\_hg\_minimum.csv}: their minima;
\item
  \texttt{data/sh\_hg\_susceptibility.csv}: the \(8\times8\)
  susceptibilities;
\item
  \texttt{data/sh\_hg\_masses.csv}: the scale, the masses and the
  amplitude mode;
\item
  \texttt{data/sh\_hg\_scan.csv}: the scan;
\item
  \texttt{data/sh\_hg\_values.csv}: single values.
\end{itemize}

The tables and numbers of this document are written from them by
\texttt{code/scripts/write\_sh\_higgs\_fragments.py}. The commands, from
the folder \texttt{Revision\_SM\_Higgs}:

\begin{Shaded}
\begin{Highlighting}[]
\CommentTok{\# the Rust program (after sh\_quantum, whose data/sh\_couplings.csv it reads)}
\KeywordTok{(}\BuiltInTok{cd}\NormalTok{ rust }\KeywordTok{\&\&} \ExtensionTok{cargo}\NormalTok{ build }\AttributeTok{{-}{-}release} \KeywordTok{\&\&} \ExtensionTok{./target/release/sh\_higgs}\KeywordTok{)}
\CommentTok{\# the Python program (after the Rust program)}
\KeywordTok{(}\BuiltInTok{cd}\NormalTok{ python }\KeywordTok{\&\&} \ExtensionTok{python3}\NormalTok{ sh\_higgs.py}\KeywordTok{)}
\CommentTok{\# the Wolfram script (on a computer with Wolfram)}
\ExtensionTok{wolframscript} \AttributeTok{{-}file}\NormalTok{ wolfram/source/sh\_higgs.wls data/sh\_hg\_minimum.csv data/sh\_hg\_masses.csv data/sh\_hg\_values.csv }\DataTypeTok{\textbackslash{}}
\NormalTok{    data/hck\_sh\_higgs\_checks\_wolfram.txt}
\CommentTok{\# the tables and numbers of this document}
\ExtensionTok{python3}\NormalTok{ code/scripts/write\_sh\_higgs\_fragments.py}
\end{Highlighting}
\end{Shaded}

The Rust program \texttt{sh\_higgs}:

\begin{longtable}[]{@{}
  >{\raggedright\arraybackslash}p{(\columnwidth - 8\tabcolsep) * \real{0.0947}}
  >{\raggedright\arraybackslash}p{(\columnwidth - 8\tabcolsep) * \real{0.0632}}
  >{\raggedright\arraybackslash}p{(\columnwidth - 8\tabcolsep) * \real{0.1053}}
  >{\raggedright\arraybackslash}p{(\columnwidth - 8\tabcolsep) * \real{0.0947}}
  >{\raggedright\arraybackslash}p{(\columnwidth - 8\tabcolsep) * \real{0.6421}}@{}}
\toprule\noalign{}
\begin{minipage}[b]{\linewidth}\raggedright
check
\end{minipage} & \begin{minipage}[b]{\linewidth}\raggedright
result
\end{minipage} & \begin{minipage}[b]{\linewidth}\raggedright
deviation
\end{minipage} & \begin{minipage}[b]{\linewidth}\raggedright
tolerance
\end{minipage} & \begin{minipage}[b]{\linewidth}\raggedright
statement
\end{minipage} \\
\midrule\noalign{}
\endhead
\bottomrule\noalign{}
\endlastfoot
H-1 & PASS & 1.778e-13 & 1.0e-9 & the state with the fixed pair
self-energy d\_0 P\_dir and a self-consistent normal self-energy is the
ground state: its energy (relative), and the pair self-energy that its
pair amplitude produces equals d\_0 P\_dir (relative; self-consistency),
five methods, both versions \\
H-2a & PASS & 9.588e-5 & 1.0e-3 & phi\_0 is a stationary point of V:
V'(phi\_0) phi\_0 / /E\_0/ from central differences with the step 0.01
in d (zero up to the error of the differences, of the order of the step
squared), five methods, both versions \\
H-2b & PASS & -- & -- & V'\,'(phi\_0) \textgreater{} 0 in every method
and both versions: phi\_0 is a minimum of the effective potential
(second HS (1, 1): 149.1883; second HS (0, 1): 143.3606; second HS (0,
0): 143.3606; second KS: 181.6833; second CJT: 181.6833; first HS (1,
1): 73.1480; first HS (0, 1): 71.6803; first HS (0, 0): 71.6803; first
KS: 80.0121; first CJT: 80.0121) \\
H-2c & PASS & 2.302e-4 & 2.0e-3 & V'\,'(phi\_0) from the steps 0.02 and
0.01 in d agree to the order of the steps (relative difference; the
stated value is their Richardson combination) \\
H-2d & PASS & 6.298e-16 & 1.0e-9 & second version: along the whole curve
the two pair amplitudes are equal, /\textless Pi\^{}d\_SH\textgreater/ =
/\textless Pi'\^{}u\_SH\textgreater/ (difference relative to phi\_0):
the curve moves both pair fields together, along the custodial
direction \\
H-3 & PASS & -- & -- & the Mexican hat along the vacuum direction: V = 0
at phi = 0 (the normal vacuum), phi grows with d, V decreases from phi =
0 to the minimum, which is at d = d\_0 among the 33 points, and
increases beyond it, five methods, both versions (second HS (1, 1): V(0)
= 5.2e-17, minimum at d = 1.00 d\_0; second HS (0, 1): V(0) = 5.2e-17,
minimum at d = 1.00 d\_0; second HS (0, 0): V(0) = 5.2e-17, minimum at d
= 1.00 d\_0; second KS: V(0) = 5.2e-17, minimum at d = 1.00 d\_0; second
CJT: V(0) = 5.2e-17, minimum at d = 1.00 d\_0; first HS (1, 1): V(0) =
5.2e-17, minimum at d = 1.00 d\_0; first HS (0, 1): V(0) = 5.2e-17,
minimum at d = 1.00 d\_0; first HS (0, 0): V(0) = 5.2e-17, minimum at d
= 1.00 d\_0; first KS: V(0) = 5.2e-17, minimum at d = 1.00 d\_0; first
CJT: V(0) = 5.2e-17, minimum at d = 1.00 d\_0) \\
H-4a & PASS & 2.562e-7 & 1.0e-6 & the susceptibilities chi\^{}\{ab\} =
d\textless Q\_a\textgreater/d mu\_b are symmetric (largest
/chi\^{}\{ab\} - chi\^{}\{ba\}/ relative to chi\^{}\{11\}), as the
theorem of Hellmann and Feynman requires (chi\^{}\{ab\} = -d\^{}2 E / d
mu\_a d mu\_b), eight generators, five methods, both versions \\
H-4b & PASS & 3.845e-4 & 1.0e-3 & the steps 0.005 and 0.01 of the
chemical potentials give the same susceptibilities to the order of the
steps squared (largest difference relative to the largest element; the
stated values are their Richardson combination) \\
H-4c & PASS & 3.941e-12 & 1.0e-6 & KS and CJT give the same
susceptibilities (relative to chi\^{}\{11\}; the perturbed states keep
their local normal matrix in the span of the Kohn-Sham densities), both
versions \\
H-4d & PASS & -- & -- & the number of broken generators, the nonzero
eigenvalues of chi (above 1e-6 of the largest; none negative): four
among the eight generators in the second version (the three of SU(2)\_L
x U(1)\_Y that the photon does not keep, and F\_tot), three among
T\^{}1, T\^{}2, T\^{}3, R, F\_tot in the first version (second HS (1,
1): 4; second HS (0, 1): 4; second HS (0, 0): 4; second KS: 4; second
CJT: 4; first HS (1, 1): 3; first HS (0, 1): 3; first HS (0, 0): 3;
first KS: 3; first CJT: 3) \\
H-4e & PASS & 1.407e-12 & 1.0e-6 & second version: the generators that
stage 3 found unbroken, T\^{}a + T\^{}a\_R (a = 1, 2, 3) and R, have
zero susceptibility with every generator (relative to chi\^{}\{11\}),
five methods \\
H-4f & PASS & 2.447e-13 & 1.0e-6 & second version: the custodial
structure chi(T\^{}a, T\^{}a) = chi(T\^{}a\_R, T\^{}a\_R) = -chi(T\^{}a,
T\^{}a\_R), equal for a = 1, 2, 3 (relative to chi\^{}\{11\}), five
methods \\
H-4g & PASS & 8.096e-13 & 1.0e-6 & the electric charge Q = T\^{}3 +
T\^{}3\_R + R/4 does not respond: chi(Q, X\_b) = 0 for the eight
generators (relative to chi\^{}\{11\}); the photon is massless, five
methods, both versions \\
H-5a & PASS & 1.703e-13 & 1.0e-6 & the mass matrix g\_a g\_b
chi\^{}\{ab\} of (W\^{}1, W\^{}2, W\^{}3, B) has one zero eigenvalue
(relative to the largest) whose eigenvector is the photon (0, 0, g',
g)/(g\^{}2 + g'\^{}2)\^{}(1/2) (one minus the squared overlap), five
methods, both versions \\
H-5b & PASS & 2.018e-13 & 1.0e-6 & its other eigenvalues: M\_W\^{}2 =
g\^{}2 chi\^{}\{11\} twice (relative), five methods, both versions \\
H-6a & PASS & 7.261e-14 & 1.0e-9 & rho from the eigenvalues, M\_W\^{}2 /
(M\_Z\^{}2 cos\^{}2 theta\_W), equals chi\^{}\{11\} / chi\^{}\{33\}
(relative), five methods, both versions \\
H-6b & PASS & 1.569e-13 & 1.0e-6 & second version: rho = chi\^{}\{11\} /
chi\^{}\{33\} = 1 (the custodial symmetry of the vacuum), five
methods \\
H-7a & PASS & -- & -- & at omega = 0 exactly four singular values of 1 -
chi(0) F vanish in the sector of relative number 0 in the second version
and three in the first (the Goldstone modes of the broken generators),
and the next is not small, HS (0, 1) and HS (0, 0) (second HS (0, 1): 4
below 1e-6, the next 2.60e-1; second HS (0, 0): 4 below 1e-6, the next
2.60e-1; first HS (0, 1): 3 below 1e-6, the next 2.60e-1; first HS (0,
0): 3 below 1e-6, the next 2.60e-1) \\
H-7b & PASS & -- & -- & below the lowest transition energy that acts on
it (the lowest omega\_c whose residue A\_c F, B\_c F on the null space
exceeds 1e-12 of the size of the bucket) there is a frequency at which 1
- chi(omega) F is singular (smallest singular value below 1e-6) and
whose null space contains the amplitude of the vacuum direction (squared
length of its projection above 0.999), HS (0, 1) and HS (0, 0), both
versions; in the first version the sector also holds the transitions of
the unpaired fields from 2 m\_LL on, whose residues on the null space
are at the level of the rounding (fraction below 1e-20): they do not act
on the mode (second HS (0, 1): omega 0.8128455463 (the lowest transition
that acts on it 0.905948, the lowest of the sector 0.905948, the largest
residue fraction below the former 0.0e0), singular value 6.6e-16,
multiplicity 4, overlap 1.000000000; second HS (0, 0): omega
0.8128455463 (the lowest transition that acts on it 0.905948, the lowest
of the sector 0.905948, the largest residue fraction below the former
0.0e0), singular value 4.8e-16, multiplicity 4, overlap 1.000000000;
first HS (0, 1): omega 0.8128455463 (the lowest transition that acts on
it 0.905948, the lowest of the sector 0.400029, the largest residue
fraction below the former 2.8e-32), singular value 5.0e-16, multiplicity
1, overlap 1.000000000; first HS (0, 0): omega 0.8128455463 (the lowest
transition that acts on it 0.905948, the lowest of the sector 0.400029,
the largest residue fraction below the former 1.3e-31), singular value
1.2e-15, multiplicity 1, overlap 1.000000000) \\
H-7c & PASS & 1.093e-15 & 1.0e-10 & HS (0, 1) and HS (0, 0) give the
same amplitude mode (relative difference of the frequencies), both
versions: the vector term of HS (0, 1), which acts on the number density
and the currents, does not couple to it \\
H-7d & PASS & 1.990e-13 & 1.0e-10 & the amplitude mode lies at the pair
self-energy, omega\_H = d\_0 (relative difference), HS (0, 1) and HS (0,
0), both versions: the analogue of the amplitude mode of the theory of
Bardeen, Cooper and Schrieffer at twice the gap parameter, here d\_0 /
2 \\
H-7e & PASS & 1.411e-5 & 1.0e-4 & the lowest transition energy that acts
on the mode is the threshold of two paired quasi-particles, 2 M\_eff =
(4 m\_LL\^{}2 + d\_0\^{}2)\^{}(1/2) (relative excess, of the order of
the smallest momentum of the grid squared; never below it), HS (0, 1)
and HS (0, 0), both versions \\
H-8a & PASS & 8.549e-14 & 1.0e-9 & the scan (HS (0, 0), 16 points in
each version): the pair self-energy d\_0 of the closed-form gap equation
(G\_SH / 4 pi\^{}2) I1((m\^{}2 + d\_0\^{}2/4)\^{}(1/2)) = 1 is
self-consistent in the engine (relative) \\
H-8b & PASS & -- & -- & in every point of the scan the amplitude mode is
found below the threshold of two quasi-particles (smallest singular
value below 1e-6, squared length of the projection of the amplitude of
the vacuum direction onto the null space above 0.999) (32 points) \\
H-8c & PASS & 9.782e-10 & 1.0e-6 & the scan point m = 0.2, G\_SH = 100.8
reproduces the main computation of HS (0, 0) in both versions: d\_0,
v\_SH and the amplitude mode (the latter from the sector of the neutral
pairs of the vacuum alone; relative) \\
H-8d & PASS & 1.030e-12 & 1.0e-9 & in every point of the scan the
amplitude mode lies at the pair self-energy, omega\_H = d\_0 (relative
difference) \\
H-9 & PASS & 1.853e-14 & 1.0e-6 & 24 and 48 radial nodes give the same
v\_SH (HS (0, 0) and CJT) and the same amplitude mode (HS (0, 0))
(relative), second version \\
\end{longtable}

The Python program \texttt{python/sh\_higgs.py}:

\begin{longtable}[]{@{}
  >{\raggedright\arraybackslash}p{(\columnwidth - 8\tabcolsep) * \real{0.0947}}
  >{\raggedright\arraybackslash}p{(\columnwidth - 8\tabcolsep) * \real{0.0632}}
  >{\raggedright\arraybackslash}p{(\columnwidth - 8\tabcolsep) * \real{0.1053}}
  >{\raggedright\arraybackslash}p{(\columnwidth - 8\tabcolsep) * \real{0.0947}}
  >{\raggedright\arraybackslash}p{(\columnwidth - 8\tabcolsep) * \real{0.6421}}@{}}
\toprule\noalign{}
\begin{minipage}[b]{\linewidth}\raggedright
check
\end{minipage} & \begin{minipage}[b]{\linewidth}\raggedright
result
\end{minipage} & \begin{minipage}[b]{\linewidth}\raggedright
deviation
\end{minipage} & \begin{minipage}[b]{\linewidth}\raggedright
tolerance
\end{minipage} & \begin{minipage}[b]{\linewidth}\raggedright
statement
\end{minipage} \\
\midrule\noalign{}
\endhead
\bottomrule\noalign{}
\endlastfoot
YH-1 & PASS & 6.455e-13 & 1.0e-09 & HS (0, 1) and HS (0, 0), both
versions: d\_0 = 0.812845546320 of the closed-form gap equation (G\_SH /
4 pi\^{}2) I1((m\^{}2 + d\_0\^{}2/4)\^{}(1/2)) = 1 against the Rust
value (relative) \\
YH-2 & PASS & 2.882e-09 & 1.0e-06 & the susceptibilities chi\^{}\{11\},
chi\^{}\{33\}, chi\^{}\{YY\}, chi\^{}\{3Y\} as -d\^{}2 E / d mu\^{}2 of
the energy of the ground state of H - mu Q (a second method) equal those
of the Rust program, d\textless Q\_a\textgreater/d mu\_b (relative to
chi\^{}\{11\}), HS (0, 0) and CJT, both versions {[}second HS (0, 0):
chi11 2.60214225e-03, chi33 2.60214225e-03, chiYY 2.60214225e-03, chi3Y
-2.60214226e-03; second CJT: chi11 3.11208804e-03, chi33 3.11208805e-03,
chiYY 3.11208805e-03, chi3Y -3.11208804e-03; first HS (0, 0): chi11
1.59943782e-03, chi33 1.30107113e-03, chiYY 1.30107113e-03, chi3Y
-1.30107113e-03; first CJT: chi11 1.88650523e-03, chi33 1.56983561e-03,
chiYY 1.56983561e-03, chi3Y -1.56983561e-03{]} \\
YH-3 & PASS & 2.016e-10 & 1.0e-04 & HS (0, 0), both versions:
V'\,'(phi\_0) from the energy and the pair amplitude of the states with
the fixed pair self-energy d P\_dir (central differences, Richardson)
against the Rust value (relative) \\
YH-4 & PASS & 6.952e-12 & 1.0e-06 & HS (0, 0): the reduced linear
response of the pair amplitudes has the phase eigenvalue 1 at omega = 0
(the Goldstone mode), no coupling of the amplitude and the phase, and
the amplitude mode below 2 M\_eff at the frequency of the Rust program;
equal in both versions (largest of these deviations) {[}second: omega
0.8128455463 (Rust 0.8128455463), threshold 0.905935, phase eigenvalue
at 0 1.000000000002, amplitude-phase coupling 1.1e-16; first: omega
0.8128455463 (Rust 0.8128455463), threshold 0.905935, phase eigenvalue
at 0 1.000000000002, amplitude-phase coupling 1.2e-16{]} \\
YH-5 & PASS & 1.768e-16 & 1.0e-09 & the masses of
data/sh\_hg\_masses.csv follow from the Rust data: M\_EFT = (V'\,'
phi\_0\^{}2 / (4 chi\^{}\{33\}))\^{}(1/2), v\_SH = 2
(chi\^{}\{11\})\^{}(1/2), rho = chi\^{}\{11\} / chi\^{}\{33\},
Lambda\_LL = v / v\_SH with v = 246.219672 GeV, M\_W = g v / 2
(relative), ten rows \\
YH-6 & PASS & 8.304e-15 & 1.0e-06 & the scan: d\_0 of the closed form
and the amplitude mode of the reduced linear response against
data/sh\_hg\_scan.csv (relative), three points {[}second m = 0.2,
G/G\^{}crit = 1.02: d\_0 0.24065630, omega 0.24065630; first m = 0.2,
G/G\^{}crit = 3: d\_0 4.01146256, omega 4.01146256; second m = 0.1,
G/G\^{}crit = 1.18384: d\_0 0.73827806, omega 0.73827806{]} \\
YH-7 & PASS & 9.853e-16 & 1.0e-12 & HS (0, 0), both versions: the
amplitude and phase eigenvalues of the reduced linear response equal the
closed forms lambda\_amp(omega) = (G/2) sum\_n w\_n (k\_n\^{}2 + m\^{}2)
/ (E\_n (E\_n\^{}2 - omega\^{}2/4)) and lambda\_phase(omega) = (G/2)
sum\_n w\_n E\_n / (E\_n\^{}2 - omega\^{}2/4), E\_n\^{}2 = k\_n\^{}2 +
m\^{}2 + d\_0\^{}2/4, and lambda\_amp(omega) - lambda\_phase(0) =
((omega\^{}2 - d\_0\^{}2)/4) S(omega) with S(omega) = (G/2) sum\_n w\_n
/ (E\_n (E\_n\^{}2 - omega\^{}2/4)) \textgreater{} 0, at five
frequencies below 2 M\_eff (absolute); lambda\_phase(0) = 1 is the gap
equation {[}second omega 0.000000: lambda\_amp 0.7377040607282 (closed
0.7377040607282), lambda\_phase 1.0000000000016 (closed
1.0000000000016); second omega 0.300000: lambda\_amp 0.7627744954223
(closed 0.7627744954223), lambda\_phase 1.0374095434543 (closed
1.0374095434543); second omega 0.600000: lambda\_amp 0.8529756584284
(closed 0.8529756584284), lambda\_phase 1.1760080331297 (closed
1.1760080331297); second omega 0.812846: lambda\_amp 1.0000000000016
(closed 1.0000000000016), lambda\_phase 1.4232626846379 (closed
1.4232626846379); second omega 0.878757: lambda\_amp 1.0858625854677
(closed 1.0858625854677), lambda\_phase 1.5946805054664 (closed
1.5946805054664); second: lambda\_phase(0) - 1 = 1.6e-12 (the gap
equation, to the tolerance of the solution); first omega 0.000000:
lambda\_amp 0.7377040607282 (closed 0.7377040607282), lambda\_phase
1.0000000000016 (closed 1.0000000000016); first omega 0.300000:
lambda\_amp 0.7627744954223 (closed 0.7627744954223), lambda\_phase
1.0374095434543 (closed 1.0374095434543); first omega 0.600000:
lambda\_amp 0.8529756584284 (closed 0.8529756584284), lambda\_phase
1.1760080331297 (closed 1.1760080331297); first omega 0.812846:
lambda\_amp 1.0000000000016 (closed 1.0000000000016), lambda\_phase
1.4232626846379 (closed 1.4232626846379); first omega 0.878757:
lambda\_amp 1.0858625854677 (closed 1.0858625854677), lambda\_phase
1.5946805054664 (closed 1.5946805054664); first: lambda\_phase(0) - 1 =
1.6e-12 (the gap equation, to the tolerance of the solution){]} \\
YH-8 & PASS & 2.442e-15 & 1.0e-09 & HS (0, 0): the number of eigenvalues
equal to 1 of the reduced linear response on the pair fields (four in
the second version at omega = 0 and at omega = d\_0; three at omega = 0
and one at omega = d\_0 in the first, on the doublet of chi); in the
second version the custodial generators T\^{}a + T\^{}a\_R commute with
the response, and the Casimir sum\_a (T\^{}a + T\^{}a\_R)\^{}2 has the
eigenvalues 0, 2, 2, 2 on the four modes at either frequency (a singlet
and a triplet; largest deviation) {[}second omega 0.000000: 4
eigenvalues equal 1, Casimir of T\^{}a + T\^{}a\_R on them
-0.0000000000, 2.0000000000, 2.0000000000, 2.0000000000; second omega
0.812846: 4 eigenvalues equal 1, Casimir of T\^{}a + T\^{}a\_R on them
-0.0000000000, 2.0000000000, 2.0000000000, 2.0000000000; first omega
0.000000: 3 eigenvalues equal 1; first omega 0.812846: 1 eigenvalues
equal 1{]} \\
YH-9 & PASS & 2.384e-15 & 1.0e-12 & the normal vacuum (no pair
self-energy) of HS (0, 0) and CJT, both versions: the ground state of H
- mu Q\_b for the eight generators and mu = +-0.005, +-0.01 has zero
charge densities and no pairs, so its susceptibilities vanish and the
gauge bosons get no mass without the condensate (largest
/\textless Q\_a\textgreater/ / /mu/ and largest element of the pair
self-energy) {[}second HS (0, 0): largest /\textless Q\_a\textgreater/ /
/mu/ 2.4e-15; second CJT: largest /\textless Q\_a\textgreater/ / /mu/
2.0e-15; first HS (0, 0): largest /\textless Q\_a\textgreater/ / /mu/
2.4e-15; first CJT: largest /\textless Q\_a\textgreater/ / /mu/
1.9e-15{]} \\
\end{longtable}

The Wolfram script \texttt{wolfram/source/sh\_higgs.wls}:

\sptkeptpenalty=10000
\begin{longtable}[]{@{}
  >{\raggedright\arraybackslash}p{(\columnwidth - 8\tabcolsep) * \real{0.0947}}
  >{\raggedright\arraybackslash}p{(\columnwidth - 8\tabcolsep) * \real{0.0632}}
  >{\raggedright\arraybackslash}p{(\columnwidth - 8\tabcolsep) * \real{0.1053}}
  >{\raggedright\arraybackslash}p{(\columnwidth - 8\tabcolsep) * \real{0.0947}}
  >{\raggedright\arraybackslash}p{(\columnwidth - 8\tabcolsep) * \real{0.6421}}@{}}
\toprule\noalign{}
\begin{minipage}[b]{\linewidth}\raggedright
check
\end{minipage} & \begin{minipage}[b]{\linewidth}\raggedright
result
\end{minipage} & \begin{minipage}[b]{\linewidth}\raggedright
deviation
\end{minipage} & \begin{minipage}[b]{\linewidth}\raggedright
tolerance
\end{minipage} & \begin{minipage}[b]{\linewidth}\raggedright
statement
\end{minipage} \\
\midrule\noalign{}
\endhead
\bottomrule\noalign{}
\endlastfoot
WH-1 & PASS & -- & -- & the effective Lagrangian Z tr((dM)\^{}dag dM) of
the bidoublet at M = phi i tau\^{}2: chi(T\^{}a, T\^{}b) =
chi(T\^{}a\_R, T\^{}b\_R) = -chi(T\^{}a, T\^{}b\_R) = Z phi\^{}2
delta\_ab; with Y = T\^{}3\_R + R/4 the charge Q = T\^{}3 + Y does not
respond; the mass matrix of (W\^{}1, W\^{}2, W\^{}3, B) has the
eigenvalues 0, g\^{}2 Z phi\^{}2 (twice), (g\^{}2 + g'\^{}2) Z phi\^{}2,
so M\_W = g v/2 and M\_Z = (g\^{}2 + g'\^{}2)\^{}(1/2) v/2 with v\^{}2 =
4 Z phi\^{}2, and rho = 1 (symbolic) (eigenvalues \{0,
g\^{}2*phi\^{}2*zz, g\^{}2*phi\^{}2*zz, (g\^{}2 +
gp\^{}2)*phi\^{}2*zz\}) \\*
WH-2 & PASS & -- & -- & the effective Lagrangian of the pair fields
gives the mass of the amplitude mode M\^{}2 = V'\,' phi\_0\^{}2 / (4
chi\^{}\{33\}) in both versions: second version chi\^{}\{33\} =
phi0\^{}2*zz, M\^{}2 = V'`/(4 Z); first version chi\^{}\{33\} =
(phi0\^{}2*zz)/2, M\^{}2 = V''/(2 Z) (symbolic) (identities) \\*
WH-3 & PASS & -- & -- & HS (0, 0): the gap equation (G\_SH / 4 pi\^{}2)
I1(M\_eff) = 1 gives M\_eff = 0.4529674056079986, d\_0 =
0.8128455463204339, phi\_0 = d\_0 I1 / (4 pi\^{}2) = d\_0 / G\_SH =
0.008063943911909066; dI1/dM = -M J with J = arsinh(1/M) - (1 +
M\^{}2)\^{}(-1/2) (exact; J equals its integral to 40 digits);
V'\,'(phi\_0) = 2 pi\^{}2 d\_0\^{}2 J / (I1 (I1 - d\_0\^{}2 J/4)) =
71.68031758679552 for one pair (first version) and twice that,
143.3606351735910, for the two pairs of the second version; against the
Rust values (relative) (deviation 1.3e-8) \\*
WH-4 & PASS & -- & -- & m\_H / v = 125.13 / v = 0.5082047224 (the inputs
of stage 2); for every row of the masses, Lambda\_LL = v / v\_SH, M\_H =
omega\_H Lambda\_LL and M\_H / v = omega\_H / v\_SH (relative); the
amplitude modes of the linear response give M\_H / v from 7.96732 to
10.1624, 15.68 to 20.00 times the value of the Standard Model (deviation
0.e-16) \\*
WH-5 & PASS & -- & -- & HS (0, 0): the identity (k\^{}2 + m\^{}2) / (E
(E\^{}2 - w\^{}2/4)) - 1/E = (w\^{}2 - d\^{}2) / (4 E (E\^{}2 -
w\^{}2/4)) with E\^{}2 = k\^{}2 + m\^{}2 + d\^{}2/4 holds for every
momentum (symbolic) and E\^{}2 - w\^{}2/4 \textgreater{} 0 for w
\textless{} 2 M\_eff (symbolic), so that lambda\_amp(w) = 1 only at w =
d\_0; the continuum integrals give /lambda\_amp(d\_0) - 1/ = 0.e-30 and
/lambda\_phase(0) - 1/ = 0.e-30 (the gap equation), the root of
lambda\_amp = 1 is 0.81284554632043386736 against d\_0 =
0.81284554632043386736, below 2 M\_eff = 0.9059348112159973 (binding
0.0930893); the amplitude modes of the Rust program (HS (0, 1) and HS
(0, 0), both versions) equal d\_0 (relative) (identity True, positivity
True; Rust deviation 8.4e-13) \\
\end{longtable}

Reference: W. A. Bardeen, C. T. Hill and M. Lindner, Minimal dynamical
symmetry breaking of the standard model, Physical Review D 41, 1647
(1990).

\hypertarget{the-symbols-of-this-document}{%
\section{9. The symbols of this
document}\label{the-symbols-of-this-document}}

\begin{longtable}[]{@{}
  >{\raggedright\arraybackslash}p{(\columnwidth - 2\tabcolsep) * \real{0.5000}}
  >{\raggedright\arraybackslash}p{(\columnwidth - 2\tabcolsep) * \real{0.5000}}@{}}
\toprule\noalign{}
\begin{minipage}[b]{\linewidth}\raggedright
symbol
\end{minipage} & \begin{minipage}[b]{\linewidth}\raggedright
meaning
\end{minipage} \\
\midrule\noalign{}
\endhead
\bottomrule\noalign{}
\endlastfoot
\(\Psi_{\rm LL}\); \(\chi_{\rm SH}\), \(\chi'_{\rm SH}\);
\(\Psi_{\rm SH}\) & the isodoublet of Dirac fields (eight components,
component \(2d+s\)), the two singlet Dirac fields (four components
each), and the three as one field of sixteen components \\
\(m_{\rm LL}\), \(m^\chi_{\rm SH}\), \(m^{\chi'}_{\rm SH}\) & the masses
of the three fields (all \(0.2\,\Lambda_{\rm LL}\)) \\
\(\Lambda_{\rm LL}\) & the momentum cutoff, the unit of the
computation \\
\(SL(2,\mathbb C)\), \(SU(2)_L\), \(U(1)_Y\) & the Lorentz group of the
spinors and the electroweak gauge group \\
\(T^{a,{\rm SM}}\), \(\tau^a\); \(Y^{\rm SM}\), \(Q^{\rm SM}\) & the
generators of \(SU(2)_L\) and the Pauli matrices; the hypercharge and
the electric charge \\
\(F^{\rm tot}_{\rm SH}\), \(R_{\rm SH}\) & the particle number and the
relative number \\
\(SU(2)_{\rm R}\), \(T^{a,{\rm R}}_{\rm SH}\); \(SU(2)_{\rm V}\) & the
custodial group, which acts on \((\chi_{\rm SH},\chi'_{\rm SH})\) as on
a doublet, and its generators; the diagonal subgroup of
\(SU(2)_L\times SU(2)_{\rm R}\) \\
\(\gamma^p_{\rm LL}\), \(\beta_{\rm LL}\), \(\alpha^j_{\rm LL}\),
\(\gamma^{\rm ch}_{\rm LL}\), \(C_{\rm LL}\); \(\sigma^{\rm LL}_j\) &
the Dirac matrices of the chiral basis, the chirality matrix, the
charge-conjugation matrix; the Pauli matrices on a Weyl spinor \\
\(\vec q_{\rm LL}\), \(q_{\rm LL}\), \(q^j_{\rm LL}\) & a momentum, its
length and its components \\
\(\mathbf e^{\rm u}_{\rm SH}\), \(\mathbf e^{\rm d}_{\rm SH}\)
(\(\mathbf e^s_{\rm SH}\)) & the isospin columns \((1,0)^T\) and
\((0,1)^T\) \\
\(\Pi^s_{\rm SH}\), \(\Pi'^s_{\rm SH}\); \(\mathsf K^s_{\rm SH}\),
\(\mathsf K'^s_{\rm SH}\) & the pair operators (1.1) and their
\(16\times16\) pair matrices \\
\(\Phi_{\rm SH}\), \(\Phi'_{\rm SH}\); \(\Phi^{\rm d}\),
\(\Phi'^{\rm u}\) & the two pair fields; their neutral components
\(\langle\Pi^{\rm d}_{\rm SH}\rangle\),
\(\langle\Pi'^{\rm u}_{\rm SH}\rangle\) \\
\(\mathcal M_{\rm SH}\) & the \(2\times2\) matrix
\((\Phi_{\rm SH},\Phi'_{\rm SH})\), a bidoublet of
\(SU(2)_L\times SU(2)_{\rm R}\) \\
\(\mathcal H_{\rm SH}\), \(\mathsf h_{\rm SH}(\vec q_{\rm LL})\) & the
Hamiltonian density (1.2) and the free Hamiltonian of the sixteen
components on a plane wave \\
\(\bar\Psi_{\rm LL}\); \({:}\ {:}\) & the Dirac adjoint; the normal
ordering with respect to the normal vacuum \\
\(G^{\rm S}_{\rm SF}\), \(G^{\rm V}_{\rm SF}\), \(G_{\rm SH}\);
\(V^{(1)}_{\rm SH}\) & the couplings of the scalar, the vector and the
pairing terms; the potential with \(G^{\rm S}_{\rm SF}=-8\),
\(G^{\rm V}_{\rm SF}=0.8\), \(G_{\rm SH}=100.8\) \\
\(G^{\rm crit}_{\rm SH}\) & the critical coupling
\(4\pi^2/I^{(1)}_{\rm SF}(m_{\rm LL})\) of the pairing terms \\
\(\mathsf N_{\rm SH}\), \(\mathsf P_{\rm SH}\);
\(\mathsf H_{\rm SH}(\vec q_{\rm LL})\) & the normal and the pair
self-energies of a uniform state; its BdG matrix (1.3) \\
\(\varrho_{\rm SH}\), \(\kappa_{\rm SH}\) & the local normal matrix,
measured from the normal vacuum, and the local pair matrix \\
\(\mathcal E^X_{\rm SH}\) & the energy per volume of a uniform state in
method \(X\), measured from the normal vacuum \\
\(X\); \(s^{\rm HS,S}_{\rm SF}\), \(s^{\rm HS,V}_{\rm SF}\) & a method
(HS (1, 1), HS (0, 1), HS (0, 0), KS or CJT); the HS parameters \\
\(\mathsf P^{\rm vac}_{\rm SH}\); \(d_{\rm SH}\), \(d^0_{\rm SH}\) & the
direction (1.4) of the vacuum in the pair self-energy; the coefficient
of a pair self-energy along it, and its value in the ground state \\
\(\varphi_{\rm SH}\), \(\varphi^0_{\rm SH}\) & the pair amplitude (1.5)
of a state, and its value in the ground state \\
\(V^{\rm eff}_{\rm SH}(\varphi_{\rm SH})\);
\(V^{\rm eff\,\prime}_{\rm SH}\), \(V^{\rm eff\,\prime\prime}_{\rm SH}\)
& the effective potential along the direction of the vacuum; its first
and second derivatives \\
\(\mathcal E\), \(\varphi\) with primes & in (2.1): derivatives of the
energy and of the pair amplitude with respect to \(d_{\rm SH}\) \\
\(M^{\rm eff}_{\rm SH}\), \(E\) & the energy
\((m_{\rm LL}^2+d_{\rm SH}^2/4)^{1/2}\) of the paired quasi-particles at
rest in the pairing-only methods, and the energy of a quasi-particle \\
\(I^{(1)}_{\rm SF}(M)\), \(J^{(3)}_{\rm SH}(M)\) & the integrals
(2.2) \\
\(M\), \(q\) & a mass and a momentum (integration variables) \\
\(Q_a\), \(Q_b\); \(\langle Q_a\rangle\) & the eight generators of the
global group (and combinations of them) as \(16\times16\) matrices; the
charge density \(\mathrm{tr}(Q_a\varrho_{\rm SH})\) \\
\(\mu^a_{\rm SH}\), \(\mu\) & the chemical potential of the charge
\(Q_a\); a chemical potential \\
\(\chi^{ab}_{\rm SH}\); \(\chi^{11}_{\rm SH}\), \(\chi^{33}_{\rm SH}\),
\(\chi^{YY}_{\rm SH}\), \(\chi^{3Y}_{\rm SH}\), \(\chi^{FF}_{\rm SH}\) &
the susceptibilities (3.1); those of \(T^{1,{\rm SM}}\), of
\(T^{3,{\rm SM}}\), of \(Y^{\rm SM}\), the mixed one of
\(T^{3,{\rm SM}}\) and \(Y^{\rm SM}\), and that of
\(F^{\rm tot}_{\rm SH}\) \\
\(W^a_\mu\), \(B_\mu\), \(A^a_0\); \(g\), \(g'\), \(g_a\); \(\theta_W\)
& the gauge fields of \(SU(2)_L\) and \(U(1)_Y\), the time component of
a gauge field; their couplings, the coupling of the generator \(a\); the
weak mixing angle \\
\(D_\mu\), \(Z_{\rm SH}\) & the covariant derivative; the kinetic
coefficient of the pair fields in the effective Lagrangian \\
\(\mathsf M^{\rm g}_{\rm SH}\) & the mass matrix (3.2) of
\((W^1_\mu,W^2_\mu,W^3_\mu,B_\mu)\) \\
\(W^\pm_\mu\), \(Z_\mu\); \(M_{W,{\rm SH}}\), \(M_{Z,{\rm SH}}\) & the
massive weak gauge fields; their masses from the superfluid \\
\(v_{\rm SH}\), \(\rho_{\rm SH}\) & the scale of the breaking (3.3) and
the ratio \(M_{W,{\rm SH}}^2/(M_{Z,{\rm SH}}^2\cos^2\theta_W)\) \\
\(v\), \(\rho\), \(m_{\rm H}\) & the vacuum value of the Higgs field of
the Standard Model, its measured parameter \(\rho\), the mass of the
Higgs boson \\
\(E^{X,{\rm gap}}_{\rm SH}\) & the gap of a ground state: the lowest
energy of one quasi-particle \\
\(\eta\) & a common oscillation of the lengths of the pair fields \\
\(M^{\rm EFT}_{\rm SH}\) & the frequency (5.1) of the amplitude mode
from the effective Lagrangian \\
\(\omega\); \(\delta\mathsf N\), \(\delta\mathsf P\),
\(\delta\bar{\mathsf P}\); \(\delta H\), \(\delta R\) & a frequency; the
amplitudes of a fluctuation of the self-energies; the changes of the BdG
matrix and of the projector onto its negative eigenvalues \\
\(E_k\), \(\lvert k\rangle\), \(f_k\) & the eigenvalues and eigenvectors
of the BdG matrix and their occupations \\
\(\boldsymbol\chi_{\rm SH}(\omega)\), \(\mathsf F^X_{\rm SH}\) & the
linear response of a uniform state at the frequency \(\omega\), and the
kernel of the contact terms of method \(X\) \\
\(\omega^{\rm H}_{\rm SH}\); \(M_{\rm H}\) & the frequency of the
amplitude mode from the linear response; the mass of the amplitude
mode \\
\(\lambda^{\rm amp}_{\rm SH}(\omega)\),
\(\lambda^{\rm ph}_{\rm SH}(\omega)\) & the eigenvalues (5.3) of the
response of the pair amplitudes in HS (0, 0) along the amplitude and
along the phase of the vacuum direction \\
\(a\), \(b\), \(d\), \(s\), \(j\), \(p\), \(k\), \(l\) & indices:
generators (or \(a=1,2,3\) of \(SU(2)_L\)), the Dirac index, the isospin
index, the space coordinates, the spacetime coordinates, eigenvectors \\
\(\mathsf 1_2\), \(\mathsf 1_4\); \({}^T\), \({}^\dagger\), \(\otimes\),
\(\mathrm{tr}\), \(\mathrm{arsinh}\) & unit matrices; the transpose, the
Hermitian conjugate, the Kronecker product, the trace, the inverse
hyperbolic sine \\
\end{longtable}

\end{document}
